% 第3章 流体动力学的基本原理 —— 依据原书扫描页手工精修版
\chapter{流体动力学的基本原理}

本章应用力学和热力学基本定律建立流体动力学的基本方程和定解条件，并根据流动的基本定律揭示流动过程的一些主要性质.作为微分型基本方程应用的特例，本章最后讲述流体静力学的主要性质.

\section*{3.1  流体动力学的积分型方程}\addcontentsline{toc}{section}{3.1 流体动力学的积分型方程}

首先讨论有限体积的流体在运动过程中的动力学方程.类似于流体运动学的两种描述方法，建立流体动力学方程也有两种方法，一种是追随流体质量体建立动力学方程；另一种是在固定的控制体内建立动力学方程.这两种分析系统的主要区别在于质量体是不变的质点系(闭系统)，而固定控制体是可变的质点系(开系统).下面先介绍这两种有限体系统的概念.

\subsection*{1. 质量体和控制体}

\noindent\textbf{定义3.1}\quad 流场中封闭流体面$\Sigma^{*}(t)$所包含的流体称为质量体.

质量体具有以下性质：

(1)质量体随流体质点一起运动，所以它的边界面形状和体积都随时间变化；

(2)质量体的边界面上没有质量输入或输出，所以质量体所包含的流体质量是不变的，它相当于热力学中的闭系统；

(3)质量体的边界面上有力的相互作用；

(4)质量体的边界面上可以有能量交换(热交换或外力做功).

例如空气中的液滴和液体中的气泡都是质量体.

\noindent\textbf{定义3.2}\quad 相对于某参照坐标系不随时间变化的封闭曲面中所包含的流体称为控制体.

控制体的性质是：

(1)控制体的几何外形和体积都是不变的；

(2)控制体的界面上可以有流体流入或流出，因此它相当于热力学中的开系统；

(3)控制体的边界面上有力的相互作用；

(4)控制体的边界面上可以有能量交换(热交换或外力做功).

例如，管道流动中由进出口限定的体积是控制体；又如，在涡轮机进出口所限定的体积也是一种控制体.

众所周知，经典的力学和热力学定理都是建立在闭系统上的，因此很容易在质量体上建立有限体的流体动力学方程.但是质量体是变形系统，质量体上的动力学方程在实际应用时很不方便，为此需要有一种方法把建立在质量体上的动力学方程变换到控制体上去，以下两节将介绍这一方法.为了明确区分这两种系统，质量体的体积和界面分别用$D^{*}(t)$，$\Sigma^{*}(t)$表示，它们是时间的函数；控制体的体积和界面分别用$D,\Sigma$表示，它们不随时间变化.

\subsection*{2. 局部导数和随体导数}

\noindent\textbf{定义3.3}\quad 控制体内某物理量总和随时间的增长率称为局部导数，用$\partial/\partial t$表示.

例如，控制体内的总质量$M=\displaystyle\int_D \rho\,\mathrm{d}V$，它的局部导数：
\begin{equation*}
\frac{\partial}{\partial t}\int_D \rho\,\mathrm{d}V=\int_D \frac{\partial\rho}{\partial t}\,\mathrm{d}V
\tag{3.1}
\end{equation*}
因为$D$是控制体包围的内域，它和时间无关，因而局部导数和控制体上的积分运算可以交换.

\noindent\textbf{定义3.4}\quad 质量体内某物理量总和对时间的增长率称为随体导数，用$D/Dt$表示.

例如，质量体内的总质量$M=\displaystyle\int_{D^{*}(t)}\rho\,\mathrm{d}V$，其随体导数为
\begin{equation*}
\frac{DM}{Dt}=\frac{D}{Dt}\int_{D^{*}(t)}\rho\,\mathrm{d}V
\tag{3.2}
\end{equation*}
因为质量体的积分域$D^{*}(t)$是随时间变化的，所以随体导数不能和质量体上的积分运算交换，但是我们有以下的输运公式.

\subsection*{3. 输运公式}

\noindent\textbf{定理3.1}\quad 任一瞬间，质量体内物理量的随体导数等于该瞬间形状、体积相同的控制体内物理量的局部导数与通过该控制体表面的输运量之和.

定理的数学表达式为
\begin{equation*}
\bigg[\frac{D}{Dt}\int_{D^{*}(t)}Q\,\mathrm{d}V\bigg]_{t=t_0}=\int_D \frac{\partial Q}{\partial t}\,\mathrm{d}V+\oint_\Sigma Q(\boldsymbol{U}\cdot\boldsymbol{n})\,\mathrm{d}A
\tag{3.3}
\end{equation*}
这里$D^{*}(t)$表示质量体，$D^{*}(t_0)$为$t=t_0$时刻的质量体，它同时等于$t=t_0$时取定的控制体体积$D$；$\Sigma$为控制体的边界面，$Q$代表任意物理量，$\displaystyle\oint_\Sigma Q(\boldsymbol{U}\cdot\boldsymbol{n})\,\mathrm{d}A$是物理量$Q$在$\Sigma$面上的输运量，又称物理量的通量.

\noindent\textbf{证明}\quad 根据随体导数的定义，左边可写成
\[
\bigg[\frac{D}{Dt}\int_{D^{*}(t)}Q\,\mathrm{d}V\bigg]_{t=t_0}=\lim_{\delta t\to 0}\frac{1}{\delta t}\bigg[\int_{D^{*}(t_0+\delta t)}Q(\boldsymbol{x},t_0+\delta t)\,\mathrm{d}V-\int_{D^{*}(t_0)}Q(\boldsymbol{x},t_0)\,\mathrm{d}V\bigg]
\]
$D^{*}(t_0+\delta t)$为时刻$t_0+\delta t$质量体的体积，我们可以把它分解为(参见图3.1)
\begin{equation*}
D^{*}(t_0+\delta t)=D^{*}(t_0)+\delta D^{*}
\tag{3.4}
\end{equation*}
$D^{*}(t_0)$是$t_0$时刻的质量体，$\delta D^{*}$是质量体$D^{*}(t_0+\delta t)$和$D^{*}(t_0)$之差.根据质量体边界面运动情况，这部分质量体的体积可以为正值，也可以是负值.

利用式(3.4)，$t_0+\delta t$时刻物理量的体积分公式可分解为两部分之和：
\[
\int_{D^{*}(t_0+\delta t)}Q(\boldsymbol{x},t_0+\delta t)\,\mathrm{d}V=\int_{D^{*}(t_0)}Q(\boldsymbol{x},t_0+\delta t)\,\mathrm{d}V+\int_{\delta D^{*}}Q(\boldsymbol{x},t_0+\delta t)\,\mathrm{d}V
\]
于是随体导数的公式可进一步写成
\begin{align*}
\bigg[\frac{D}{Dt}\int_{D^{*}(t)}Q\,\mathrm{d}V\bigg]&=\lim_{\delta t\to 0}\frac{1}{\delta t}\bigg[\int_{D^{*}(t_0)}Q(\boldsymbol{x},t_0+\delta t)\,\mathrm{d}V-\int_{D^{*}(t_0)}Q(\boldsymbol{x},t_0)\,\mathrm{d}V\bigg]\\
&\quad+\lim_{\delta t\to 0}\frac{1}{\delta t}\bigg[\int_{\delta D^{*}}Q(\boldsymbol{x},t_0+\delta t)\,\mathrm{d}V\bigg]
\end{align*}
上式右边第一项是在指定域$D^{*}(t_0)=D$上的局部导数，故该项可写成
\[
\lim_{\delta t\to 0}\frac{1}{\delta t}\bigg[\int_{D^{*}(t_0)}Q(\boldsymbol{x},t_0+\delta t)\,\mathrm{d}V-\int_{D^{*}(t_0)}Q(\boldsymbol{x},t_0)\,\mathrm{d}V\bigg]=\int_{D^{*}(t_0)}\frac{\partial Q}{\partial t}\,\mathrm{d}V=\int_D \frac{\partial Q}{\partial t}\,\mathrm{d}V
\]
前式右边第2项可以化为面积分.具体做法如下：将$\delta D^{*}$无限分割为以$\Sigma(t_0)$面上的微元面积为底，以$\Sigma(t_0)$上任意质点到$\Sigma(t_0+\delta t)$面的位移为边的无穷多微元斜柱体.由$\Sigma(t_0)$上任意质点到$\Sigma(t_0+\delta t)$面的位移等于$\boldsymbol{U}\delta t$，因而斜柱体的体积应为
\[
\mathrm{d}V=(\boldsymbol{U}\delta t)\cdot\boldsymbol{n}\,\mathrm{d}A
\]
$\boldsymbol{n}$为微元面积的法向量.应当注意到，如果$\boldsymbol{U}\cdot\boldsymbol{n}>0$，这时$\mathrm{d}V>0$，也就是$\delta D^{*}$中的正值的部分；反之$\boldsymbol{U}\cdot\boldsymbol{n}<0$，则$\mathrm{d}V<0$，它是原质量体的减小部分.将$\mathrm{d}V=(\boldsymbol{U}\delta t)\cdot\boldsymbol{n}\,\mathrm{d}A$代入前面随体导数的体积分式后，在$\delta D^{*}$域上体积分可写成
\begin{align*}
\lim_{\delta t\to 0}\bigg[\frac{1}{\delta t}\int_{\delta D^{*}}Q(\boldsymbol{x},t_0+\delta t)\,\mathrm{d}V\bigg]&=\lim_{\delta t\to 0}\bigg[\frac{1}{\delta t}\oint_{\Sigma^{*}(t_0)}Q(\boldsymbol{x},t_0+\delta t)(\boldsymbol{U}\cdot\boldsymbol{n})\,\mathrm{d}A\,\delta t\bigg]\\
&=\oint_{\Sigma^{*}(t_0)}Q(\boldsymbol{x},t_0)(\boldsymbol{U}\cdot\boldsymbol{n})\,\mathrm{d}A=\oint_\Sigma Q(\boldsymbol{x},t_0)(\boldsymbol{U}\cdot\boldsymbol{n})\,\mathrm{d}A
\end{align*}
这一项等于物理量$Q$在控制体界面上的输运量，即由流体运动携带出去的量.我们约定指向控制体外的法向量为正，就是说正值的输运量$(\boldsymbol{U}\cdot\boldsymbol{n}>0)$表示由控制体输出的物理量，负值的输运量$(\boldsymbol{U}\cdot\boldsymbol{n}<0)$表示流入控制体的物理量，于是证明了质量体上随体导数的公式(3.3)，该式可概括表述为

\textbf{随体导数$=$局部导数$+$控制体输出的输运量}

下面先建立质量体上的动力学守恒方程，然后应用输运公式建立控制体上的动力学守恒方程.

\subsection*{4. 质量体上的动力学方程}

应用力学和热力学定律，质量体上有以下基本方程.

(1)质量守恒方程

质量体是闭系统，因而它的总质量不变，即
\begin{equation*}
\frac{D}{Dt}\int_{D^{*}(t)}\rho\,\mathrm{d}V=0
\tag{3.5}
\end{equation*}

(2)动量方程

根据牛顿定律，质量体内动量增长率等于该瞬间作用在质量体上的外力之和.将牛顿定律应用于流体的质量体，就有质量体上的动量方程：

质量体内动量增长率为
\[
\frac{D}{Dt}\int_{D^{*}(t)}\rho\boldsymbol{U}\,\mathrm{d}V
\]
质量体内体积力之和为
\[
\int_{D^{*}(t)}\rho\boldsymbol{f}\,\mathrm{d}V
\]
质量体边界面上表面力之和为
\[
\oint_{\Sigma^{*}(t)}\boldsymbol{T}_n\,\mathrm{d}A
\]
于是动量方程为
\begin{equation*}
\frac{D}{Dt}\int_{D^{*}(t)}\rho\boldsymbol{U}\,\mathrm{d}V=\int_{D^{*}(t)}\rho\boldsymbol{f}\,\mathrm{d}V+\oint_{\Sigma^{*}(t)}\boldsymbol{T}_n\,\mathrm{d}A
\tag{3.6}
\end{equation*}

(3)动量矩方程

质点系动量矩定理可陈述如下：质点系的动量矩增长率等于该瞬间作用在质点系上外力矩之和.应用质点系动量矩定理，类似式(3.6)的导出过程，可以导出质量体上的动量矩方程为
\begin{equation*}
\frac{D}{Dt}\int_{D^{*}(t)}\rho\boldsymbol{r}\times\boldsymbol{U}\,\mathrm{d}V=\int_{D^{*}(t)}\rho\boldsymbol{r}\times\boldsymbol{f}\,\mathrm{d}V+\oint_{\Sigma^{*}(t)}\boldsymbol{r}\times\boldsymbol{T}_n\,\mathrm{d}A
\tag{3.7}
\end{equation*}
式中$\boldsymbol{r}$为流场中任意点到参考点的向径，上式左边是质量体中总动量矩的时间增长率，右边第一项是作用在质量体内体积力的合力矩，右边第二项是作用在质量体表面上的表面力的合力矩.

(4)能量守恒方程

遵照热力学第一定律，质量体(闭系统)内总能量的增长率等于单位时间内的外力做功和输入质量体内的热量之和.

质量体的总能量等于动能和内能之和，故质量体内的总能量增长率应为
\[
\frac{D}{Dt}\int_{D^{*}(t)}\rho\Big(e+\frac{|\boldsymbol{U}|^{2}}{2}\Big)\mathrm{d}V
\]
$e$为流体单位质量的内能，它是热力学状态参数；$|\boldsymbol{U}|^{2}/2$是单位质量的动能.外力做功和输入的热量有

体积力所做功率：$\displaystyle\int_{D^{*}(t)}\rho\boldsymbol{f}\cdot\boldsymbol{U}\,\mathrm{d}V$

表面力所做功率：$\displaystyle\oint_{\Sigma^{*}(t)}\boldsymbol{T}_n\cdot\boldsymbol{U}\,\mathrm{d}A$

质量体内的生成热：$\displaystyle\int_{D^{*}(t)}\rho\dot{q}\,\mathrm{d}V$

边界面上因热传导输入的热量：$\displaystyle\oint_{\Sigma^{*}(t)}\lambda\boldsymbol{n}\cdot\mathrm{grad}\,T\,\mathrm{d}A$

于是质量体的能量守恒方程为
\begin{align*}
\frac{D}{Dt}\int_{D^{*}(t)}\rho\Big(e+\frac{|\boldsymbol{U}|^{2}}{2}\Big)\mathrm{d}V&=\int_{D^{*}(t)}\rho\boldsymbol{f}\cdot\boldsymbol{U}\,\mathrm{d}V+\oint_{\Sigma^{*}(t)}\boldsymbol{T}_n\cdot\boldsymbol{U}\,\mathrm{d}A\\
&\quad+\int_{D^{*}(t)}\rho\dot{q}\,\mathrm{d}V+\oint_{\Sigma^{*}(t)}\lambda\boldsymbol{n}\cdot\mathrm{grad}\,T\,\mathrm{d}A
\tag{3.8}
\end{align*}
式中$\dot{q}$是流体微团单位质量在单位时间内的生成热，例如：有化学反应时的化学反应热等；$\lambda$是傅里叶导热系数.

守恒方程组(3.5)$\sim$(3.8)的左边都是随体导数，利用输运公式(3.3)就可以把它们转化为控制体上的守恒方程，在实际使用时就方便得多了.

\subsection*{5. 控制体上的守恒方程}

利用式(3.3)，用$t_0$时刻的局部导数及输运量来表示该时刻的随体导数.

质量体内的质量增长率
\begin{equation*}
\frac{D}{Dt}\int_{D^{*}}\rho\,\mathrm{d}V=\int_{D^{*}(t)}\frac{\partial\rho}{\partial t}\,\mathrm{d}V+\oint_{\Sigma^{*}(t)}\rho(\boldsymbol{U}\cdot\boldsymbol{n})\,\mathrm{d}A
\tag{3.9}
\end{equation*}
质量体上的动量增长率
\begin{equation*}
\frac{D}{Dt}\int_{D^{*}}\rho\boldsymbol{U}\,\mathrm{d}V=\int_{D^{*}(t)}\frac{\partial\rho\boldsymbol{U}}{\partial t}\,\mathrm{d}V+\oint_{\Sigma^{*}(t)}\rho\boldsymbol{U}(\boldsymbol{U}\cdot\boldsymbol{n})\,\mathrm{d}A
\tag{3.10}
\end{equation*}
质量体上的动量矩增长率
\begin{equation*}
\frac{D}{Dt}\int_{D^{*}}\rho\boldsymbol{r}\times\boldsymbol{U}\,\mathrm{d}V=\int_{D^{*}(t)}\frac{\partial}{\partial t}(\rho\boldsymbol{r}\times\boldsymbol{U})\,\mathrm{d}V+\oint_{\Sigma^{*}(t)}\rho\boldsymbol{r}\times\boldsymbol{U}(\boldsymbol{U}\cdot\boldsymbol{n})\,\mathrm{d}A
\tag{3.11}
\end{equation*}
质量体上的能量增长率
\begin{align*}
\frac{D}{Dt}\int_{D^{*}}\rho\Big(e+\frac{|\boldsymbol{U}|^{2}}{2}\Big)\mathrm{d}V&=\int_{D^{*}(t)}\frac{\partial}{\partial t}\Big[\rho\Big(e+\frac{|\boldsymbol{U}|^{2}}{2}\Big)\Big]\mathrm{d}V\\
&\quad+\oint_{\Sigma^{*}(t)}\rho\Big(e+\frac{|\boldsymbol{U}|^{2}}{2}\Big)(\boldsymbol{U}\cdot\boldsymbol{n})\,\mathrm{d}A
\tag{3.12}
\end{align*}
将式(3.9)$\sim$式(3.12)代入式(3.5)$\sim$式(3.8)，并取控制体$D=D^{*}(t)$，$\Sigma=\Sigma^{*}(t)$得到下列控制体上的守恒方程.

(1)质量守恒方程
\begin{equation*}
\int_D \frac{\partial\rho}{\partial t}\,\mathrm{d}V+\oint_\Sigma \rho(\boldsymbol{U}\cdot\boldsymbol{n})\,\mathrm{d}A=0
\tag{3.13}
\end{equation*}

(2)动量方程
\begin{equation*}
\int_D \frac{\partial(\rho\boldsymbol{U})}{\partial t}\,\mathrm{d}V+\oint_\Sigma \rho\boldsymbol{U}(\boldsymbol{U}\cdot\boldsymbol{n})\,\mathrm{d}A=\int_D \rho\boldsymbol{f}\,\mathrm{d}V+\oint_\Sigma \boldsymbol{T}_n\,\mathrm{d}A
\tag{3.14}
\end{equation*}

(3)动量矩方程
\begin{equation*}
\int_D \frac{\partial(\rho\boldsymbol{r}\times\boldsymbol{U})}{\partial t}\,\mathrm{d}V+\oint_\Sigma \rho\boldsymbol{r}\times\boldsymbol{U}(\boldsymbol{U}\cdot\boldsymbol{n})\,\mathrm{d}A=\int_D \rho\boldsymbol{r}\times\boldsymbol{f}\,\mathrm{d}V+\oint_\Sigma \boldsymbol{r}\times\boldsymbol{T}_n\,\mathrm{d}A
\tag{3.15}
\end{equation*}

(4)能量守恒方程
\begin{align*}
\int_D \frac{\partial}{\partial t}\Big[\rho\Big(e+\frac{|\boldsymbol{U}|^{2}}{2}\Big)\Big]\mathrm{d}V&+\oint_\Sigma \rho\Big(e+\frac{|\boldsymbol{U}|^{2}}{2}\Big)(\boldsymbol{U}\cdot\boldsymbol{n})\,\mathrm{d}A\\
&=\int_D \rho\boldsymbol{f}\cdot\boldsymbol{U}\,\mathrm{d}V+\oint_\Sigma \boldsymbol{T}_n\cdot\boldsymbol{U}\,\mathrm{d}A+\int_D \rho\dot{q}\,\mathrm{d}V+\oint_\Sigma \lambda\,\frac{\partial T}{\partial n}\,\mathrm{d}A
\tag{3.16}
\end{align*}

\subsection*{6. 非惯性坐标系中的守恒方程}

控制体方法通常应用于流体的定常流动，但是在不同参照坐标系统中流动的定常性是不同的，例如地球上的定常流动，在太阳系的参照坐标系中观察时是非定常的；又如：旋转的流体机械内部流动，对于固定在旋转叶轮的参照坐标系中考察是定常的；而从地面参照坐标来考察该流动是非定常的.为了便于应用控制体守恒方程，如果在运动参照坐标系中流动是定常的，往往在运动参照坐标系中建立控制体，这时就需要导出运动参照坐标系中的守恒型方程.

根据质点运动学和动力学原理，在运动坐标系中质点相对运动的动力学方程中应当包含惯性力强度$-\boldsymbol{a}$，它是相对于绝对坐标系的质点牵连加速度及质点科氏加速度之和的负值.质点牵连加速度就是运动坐标系中对应点相对于绝对坐标系的加速度；科氏加速度等于质点相对速度和运动坐标系的旋转角速度的叉积.即$\boldsymbol{a}$的一般形式为
\begin{equation*}
\boldsymbol{a}=\dot{\boldsymbol{V}}_0(t)+\dot{\boldsymbol{\omega}}(t)\times\boldsymbol{r}+\boldsymbol{\omega}\times(\boldsymbol{\omega}\times\boldsymbol{r})+2\boldsymbol{\omega}\times\boldsymbol{U}
\tag{3.17}
\end{equation*}
式中$\dot{\boldsymbol{V}}_0(t)$是运动坐标系原点的移动加速度，$\boldsymbol{\omega}$是运动坐标系相对于绝对坐标系的角速度，$\dot{\boldsymbol{\omega}}(t)$是角加速度.因此$\dot{\boldsymbol{\omega}}(t)\times\boldsymbol{r}$是切向加速度，$\boldsymbol{\omega}\times(\boldsymbol{\omega}\times\boldsymbol{r})$是向心加速度，式(3.17)右边前3项之和为牵连加速度；$\boldsymbol{U}$是运动坐标系中当地流体质点的相对运动速度，$2\boldsymbol{\omega}\times\boldsymbol{U}$是科氏加速度.在控制体守恒方程(3.14)$\sim$(3.16)中将$\boldsymbol{f}$换成$\boldsymbol{f}-\boldsymbol{a}$就能得到非惯性坐标系中的积分型守恒方程.后面3.2节中将用具体例题说明如何应用上述公式.

\subsection*{7. 理想流体管内一维流动的常用公式}

流体的管内流动是最简单而又常见的流动.在近似分析计算中，假定管内流动是单向、均匀的，即在管截面上流动速度是相等的，并都沿同一方向运动.这种近似称为一维流动假定，这时管截面$A$是管流轴向坐标$x$的函数，截面上的流速$U$，压强$p$，密度$\rho$等物理量也只是轴向坐标$x$的函数.下面应用控制体上理想流体流动的守恒方程导出几个常用的公式，具体应用在3.2节中介绍.

任取长度$\mathrm{d}x$的管流，以它的入口、出口和管壁为控制体，如图3.2所示.在该控制体上建立质量动量和能量的守恒方程.

(1)管内定常流的质量守恒方程

定常流中局部导数为零，任意管截面上的质量流量相等，故质量守恒方程为
\begin{equation*}
\oint_\Sigma \rho(\boldsymbol{U}\cdot\boldsymbol{n})\,\mathrm{d}A=0
\tag{3.18}
\end{equation*}
这里积分面积$\Sigma$包括：管的侧面$\delta A_{\text{管壁}}$和进出口端面$A_1,A_2$(图3.2)，于是式(3.18)可写作
\[
\oint_{A_{\text{管壁}}}\rho(\boldsymbol{U}\cdot\boldsymbol{n})\,\mathrm{d}A+\oint_{A_1}\rho(\boldsymbol{U}\cdot\boldsymbol{n})\,\mathrm{d}A+\oint_{A_2}\rho(\boldsymbol{U}\cdot\boldsymbol{n})\,\mathrm{d}A=0
\]
管壁侧面是不可穿透的，在该面上有$\boldsymbol{U}\cdot\boldsymbol{n}=0$，故通过流管的质量守恒方程为
\begin{equation*}
\oint_{A_2}\rho(\boldsymbol{U}\cdot\boldsymbol{n})\,\mathrm{d}A=-\oint_{A_1}\rho(\boldsymbol{U}\cdot\boldsymbol{n})\,\mathrm{d}A
\tag{3.19a}
\end{equation*}
注意到$\boldsymbol{n}$总是指向控制体外，故式(3.19a)表明：定常流由截面$A_1$流入(式中负号表示流入)的质量流量等于由截面$A_2$流出的质量流量.

对于均匀不可压缩流体，即$\rho=\mathrm{const.}$，可进一步导出
\begin{equation*}
\oint_{A_2}(\boldsymbol{U}\cdot\boldsymbol{n})\,\mathrm{d}A=-\oint_{A_1}(\boldsymbol{U}\cdot\boldsymbol{n})\,\mathrm{d}A
\tag{3.19b}
\end{equation*}
这里积分式表示单位时间通过管截面的流体体积，即体积流量.因此，在不可压缩流体的管内流动中，任意管截面上体积流量相等.

在一维近似中可压缩流体定常管内流动的连续方程为
\begin{equation*}
\rho_1 U_1 A_1=\rho_2 U_2 A_2=\rho(x)U(x)A(x)
\tag{3.19c}
\end{equation*}
不可压缩流体管内流动的连续方程为
\begin{equation*}
U_1 A_1=U_2 A_2=U(x)A(x)
\tag{3.19d}
\end{equation*}
对于不可压缩流体的管内流动，控制体中的质量和体积是不变的，即控制体内质量的局部导数始终等于零，因此式(3.19d)也可以用于非定常不可压缩流体的管内流动.不可压缩流体管内流动的连续方程的微分形式可写作
\begin{equation*}
\mathrm{d}[U(x)A(x)]=0
\tag{3.19e}
\end{equation*}

(2)理想不可压缩流体在管内流动的动量守恒方程

由于管壁是不可穿透的，因此控制体内动量增长率包括：通过控制体的动量通量和单位时间控制体内的动量增量.单位时间通过控制体的动量通量等于$(U_2U_2A_2-U_1U_1A_1)$，利用不可压缩流体的连续方程$U_1A_1=U_2A_2$，通过控制体的动量通量可表示为

\textbf{单位时间通过控制体的动量通量}$=(U_2U_2A_2-U_1U_1A_1)=AU\mathrm{d}U$

控制体内的动量增长率很容易求出

\textbf{单位时间控制体内的动量增量}$=\partial(\rho U\delta V)/\partial t$

作用在控制体上的外力包括：压强合力和体积力的合力.

\textbf{控制体上压强的合力}$=p_1A_1-p_2A_2-(p_1+p_2)\delta A_{\text{管壁}}\,n_x/2$

$n_x$是管壁面的单位法向量在$x$方向的投影，由图3.2容易导出$\delta A_{\text{管壁}}\,n_x=(A_1-A_2)=-\mathrm{d}A$，因此$(p_1+p_2)\delta A_{\text{管壁}}\,n_x/2=\bar{p}\,\mathrm{d}A+$高阶小量，$\bar{p}=(p_1+p_2)/2$，于是

\textbf{控制体上压强的合力}$=p_1A_1-p_2A_2+\bar{p}\,\mathrm{d}A=-\mathrm{d}(pA)+\bar{p}\,\mathrm{d}A=-A\,\mathrm{d}p$

控制体上体积力的合力很容易求出

\textbf{控制体上体积力的合力}$=\rho f_x\delta V$

根据控制体上的动量守恒方程：\textbf{控制体上的动量增长率和单位时间动量通量之和等于作用在控制体上的外力之和}，于是有
\[
\partial(\rho U\delta V)/\partial t+\rho A U\,\mathrm{d}U=-A\,\mathrm{d}p+\rho f_x\,\delta V
\]
控制体的体积$\delta V=A\mathrm{d}x+$高阶小量，上式可简化为
\begin{equation*}
\frac{\partial U}{\partial t}+U\,\frac{\mathrm{d}U}{\mathrm{d}x}=-\frac{1}{\rho}\,\frac{\mathrm{d}p}{\mathrm{d}x}+f_x
\tag{3.20}
\end{equation*}
在有势力场中，$f_x=-\partial\Pi/\partial x$，上式可简化为
\begin{equation*}
\frac{\partial U}{\partial t}+\frac{\mathrm{d}}{\mathrm{d}x}\Big(\frac{U^{2}}{2}+\frac{p}{\rho}+\Pi\Big)=0
\tag{3.21a}
\end{equation*}
不可压缩流体作定常$(\partial U/\partial t=0)$运动时，上式可以得到简单的积分公式
\begin{equation*}
\frac{p}{\rho}+\frac{U^{2}}{2}+\Pi=\mathrm{const.}
\tag{3.21b}
\end{equation*}
式(3.21b)常称为伯努利公式，或伯努利积分.

由于定常流动的流管也是不可穿透的，因此式(3.21b)也适用于不可压缩理想流体定常运动的流管.

(3)理想流体在势力场中作定常绝热管内流动时的能量方程

应用控制体上的能量守恒方程可以导出以下理想流体在势力场中作定常绝热管内流动时的能量方程
\begin{equation*}
\frac{p}{\rho}+e+\frac{U^{2}}{2}+\Pi=C
\tag{3.22a}
\end{equation*}
控制体上的能量守恒方程为
\[
\oint_\Sigma \rho\Big(e+\frac{|\boldsymbol{U}|^{2}}{2}\Big)(\boldsymbol{U}\cdot\boldsymbol{n})\,\mathrm{d}A=-\int_D \rho(\mathrm{grad}\,\Pi)\cdot\boldsymbol{U}\,\mathrm{d}V-\oint_\Sigma p(\boldsymbol{n}\cdot\boldsymbol{U})\,\mathrm{d}A+\int_D \rho\dot{q}\,\mathrm{d}V
\]
将能量方程应用到管流的控制体上，如图3.2所示，在管侧面$\boldsymbol{U}\cdot\boldsymbol{n}=0$，因此能量输运量简化为
\[
\oint_\Sigma \rho\Big(e+\frac{|\boldsymbol{U}|^{2}}{2}\Big)(\boldsymbol{U}\cdot\boldsymbol{n})\,\mathrm{d}A=\rho_2\Big(e_2+\frac{U_2^{2}}{2}\Big)U_2A_2-\rho_1\Big(e_1+\frac{U_1^{2}}{2}\Big)U_1A_1
\]
应用微元流管上的质量守恒方程(3.19a)：$\rho_1U_1A_1=\rho_2U_2A_2$，上式可简化为
\[
\oint_\Sigma \rho\Big(e+\frac{|\boldsymbol{U}|^{2}}{2}\Big)(\boldsymbol{U}\cdot\boldsymbol{n})\,\mathrm{d}A=\bigg[\Big(e_2+\frac{U_2^{2}}{2}\Big)-\Big(e_1+\frac{U_1^{2}}{2}\Big)\bigg]\rho_1U_1A_1
\]
对于压强做功项，也做同样的推导：
\[
-\oint_\Sigma p(\boldsymbol{U}\cdot\boldsymbol{n})\,\mathrm{d}A=p_1U_1A_1-p_2U_2A_2=\Big(\frac{p_1}{\rho_1}-\frac{p_2}{\rho_2}\Big)\rho_1U_1A_1
\]
在势力场做功项$-\displaystyle\int\rho(\mathrm{grad}\,\Pi)\cdot\boldsymbol{U}\,\mathrm{d}V$中，令体积元$\mathrm{d}V=A\mathrm{d}x$，于是：
\begin{align*}
-\int_1^2 \rho(\mathrm{grad}\,\Pi)\cdot\boldsymbol{U}\,\mathrm{d}V&=-\int_1^2 \rho(\mathrm{grad}\,\Pi)\cdot\boldsymbol{U}A\,\mathrm{d}x\\
&=-\int_1^2(\rho U A)\,\frac{\partial\Pi}{\partial x}\,\mathrm{d}x=-\rho U A(\Pi_2-\Pi_1)
\end{align*}
在绝热流动中$\dot{q}=0$，将上面导出的能量方程的各项代回原方程，并应用质量守恒方程消去$\rho U A$，得沿管流的能量方程为
\[
\frac{p}{\rho}+e+\frac{U^{2}}{2}+\Pi=C
\]
于是证明了式(3.22a).在热力学中$e+p/\rho=h$是单位质量的热焓，因此式(3.22a)的意义是：

\textbf{理想流体的定常绝热管内流动中单位质量的焓、动能和势能之和是常数.}

对于理想的完全气体，热焓等于比定压热容和绝对温度的乘积，即
\[
h=c_p T
\]
在忽略势能的条件下，理想完全气体的绝热管内定常流动的能量方程为
\begin{equation*}
c_p T+\frac{U^{2}}{2}=\mathrm{const.}
\tag{3.22b}
\end{equation*}
上式表明，理想完全气体的定常绝热流动中，速度增加，温度降低.由于定常流体运动的流管也是不可穿透的，因此式(3.22a)，式(3.22b)也可以应用于理想完全气体定常绝热流动的流管.

\section*{3.2  伯努利公式的应用}\addcontentsline{toc}{section}{3.2 伯努利公式的应用}

(1)理想不可压缩流体定常管内流动

前面已经导出不可压缩理想流体的管内定常流动的伯努利公式：
\[
\frac{U^{2}}{2}+\frac{p}{\rho}+\Pi=\mathrm{const.}
\]
它的物理意义简单明了：在没有外力的情况下，理想(无粘性)不可压缩流体的管内流动中，速度大的截面处，压强低；速度小的截面处，压强高.还可以从能量守恒的角度理解伯努利公式，说明单位质量流体的动能(第一项)、势能(第三项)与压强除以密度之和在管道流动中是守恒的.伯努利公式在工程计算中非常有用，下面举一些简单的例子予以说明.

\noindent\textbf{例3.1}\quad 盛液容器中小孔排液的计算.

设有盛液的巨大容器(如水库或储液罐)，在液面下容器底部有一排液小孔，假定液体粘性可以忽略，已知液面上压强为$p_1$，孔口外压强为$p_2$，孔口面积为$a$，计算小孔泻出的流量(图3.3).

\noindent\textbf{解}\quad 在液面和孔口间利用不可压缩流体流动的连续性方程，则有
\[
U_1 A=U_2 a
\]
$A$为容器中液面面积，因盛液容器很大，$a/A\ll 1$，因而$U_1\ll U_2$可以忽略不计，这表明容器液面几乎保持不变，因而可以把流动近似为定常的，应用定常理想不可压缩管内流动的伯努利公式：
\[
\frac{p_1}{\rho_1}+\frac{U_1^{2}}{2}+\Pi_1=\frac{p_2}{\rho_2}+\frac{U_2^{2}}{2}+\Pi_2
\]
单位质量的势能差$\Pi_1-\Pi_2=gH$，又有$U_1\ll U_2$，故
\[
\frac{U_2^{2}}{2}=\frac{p_2-p_1}{\rho}+gH
\]
\[
U_2=\sqrt{2(p_2-p_1)/\rho+2gH}
\]
孔口外压强和液面压强相等，$p_1=p_2$，于是
\[
U_2=\sqrt{2gH}
\]
上式表明，理想不可压缩流体在重力场作用下，在水面下$H$的小孔出流的速度等于同样高度下物体在真空中自由落体的速度.小孔排出的体积流量为
\[
Q=a\sqrt{2gH}
\]

\noindent\textbf{例3.2}\quad 文丘利(Venturi)流量计.

这是连接在管路中测量不可压缩流体定常流动体积流量的一种常用仪器(图3.4)，当测定流动的进口与喉部压差和已知进口与喉部面积时，便可以计算通过文丘利管的流量.

\noindent\textbf{解}\quad 流动是定常、理想和不可压缩的，假定在管流截面上流速均匀分布.首先利用不可压缩流体管流的连续方程
\[
U_1 A_1=U_2 A_2
\]
另外有伯努利公式
\[
\frac{p_1}{\rho_1}+\frac{U_1^{2}}{2}+\Pi_1=\frac{p_2}{\rho_2}+\frac{U_2^{2}}{2}+\Pi_2
\]
设文丘利管为水平放置，则$\Pi_1=\Pi_2$，于是有
\[
\frac{U_2^{2}-U_1^{2}}{2}=\frac{p_1-p_2}{\rho}
\]
将连续方程代入伯努利方程后有
\[
\frac{U_2^{2}}{2}-\Big(\frac{A_2}{A_1}\Big)^{2}\frac{U_2^{2}}{2}=\frac{p_1-p_2}{\rho}
\]
用U形管测量压差$p_1-p_2$，则
\[
p_1-p_2=(\rho'-\rho)gH
\]
式中$\rho$和$\rho'$分别为U形管两臂中液体的密度，于是
\[
U_2=\sqrt{\frac{2(\rho'-\rho)gH}{\rho\Big[1-\Big(\dfrac{A_2}{A_1}\Big)^{2}\Big]}}
\]
体积流量为
\[
Q=A_2U_2=A_2\sqrt{\frac{2(\rho'-\rho)gH}{\rho\Big[1-\Big(\dfrac{A_2}{A_1}\Big)^{2}\Big]}}
\]
上式是理想文丘利管的流量计算公式.真实流体是有粘性的，并且流动在截面上是近似均匀的，因此实际文丘利管的流量公式需要修正，常用工程计算公式为
\[
Q=\eta A_2\sqrt{\frac{2(\rho'-\rho)gH}{\rho\Big[1-\Big(\dfrac{A_2}{A_1}\Big)^{2}\Big]}}
\]
式中$\eta$为修正系数.设计良好的文丘利管的流量系数$\eta>0.90$，作为工程应用来说，理想流体公式是相当好的.

\noindent\textbf{例3.3}\quad 利用水银U形管压强计，测得输水管路中文丘利管压差为$H=5$cm汞柱，并已知$A_2/A_1=1/4$，$A_2=50\,\mathrm{cm^2}$，计算通过文丘利管的喉部流速和流量，已知$\eta=0.92$.

\noindent\textbf{解}\quad 利用例3.2的公式，已知$\rho'/\rho=13.6$
\begin{gather*}
U_2=\sqrt{\frac{2\times 12.6\times 9.81\times 0.05}{1-(0.25)^{2}}}=3.63\,(\mathrm{m/s})\\[4pt]
Q=\eta A_2U_2=0.92\times 0.0050\times U_2=0.0167\,(\mathrm{m^3/s})
\end{gather*}

(2)理想不可压缩流体的管内非定常流动

\noindent\textbf{例3.4}\quad U形管中液面的震荡(图3.5).

有一支垂直放置的等截面U形管内盛有不可压缩液体，由于晃动，U形管两臂中初始液有高差$H$.由于液面高差，管内液体失去平衡，发生震荡，假定流体是理想的，U形管液柱的总长度是$L$，计算震荡的频率和液面的运动规律.

\noindent\textbf{解}\quad 在等截面U形管中，根据连续方程，速度$U$在液柱内是常数.利用式(3.21a)，沿液柱从高液面到低液面间积分，得
\[
\int_0^{L}\frac{\partial U}{\partial t}\,\mathrm{d}x+\int_0^{L}\frac{\partial(p/\rho+\Pi)}{\partial x}\,\mathrm{d}x=0
\]
U形管两边液面上压强都是大气压强，因此压强项的积分等于零.积分式的结果是
\[
\frac{\partial U}{\partial t}\,L+\Pi(L)-\Pi(0)=0
\]
$\Pi(L)-\Pi(0)=-gh$，故
\[
\frac{\partial U}{\partial t}\,L-gh=0
\]
以平衡液面作为坐标面，左液面坐标是$h/2$，右液面坐标是$-h/2$.震荡时，左液面下降，故
\[
U=-\frac{1}{2}\,\frac{\mathrm{d}h}{\mathrm{d}t}
\]
代入前面公式，得
\[
\frac{\partial^{2}h}{\partial t^{2}}+\frac{2g}{L}\,h=0
\]
以上常微分方程的解是
\[
h=A\sin\omega t+B\cos\omega t
\]
$\omega^{2}=2g/L$，震荡频率$\omega=\sqrt{2g/L}$.初始时刻$t=0$：液柱处于静止状态，$h=H$，$\dfrac{\mathrm{d}h}{\mathrm{d}t}=U=0$，于是液面的运动方程为
\[
h=H\cos\omega t
\]

\section*{3.3  定常流控制体积分型守恒方程的应用}\addcontentsline{toc}{section}{3.3 定常流控制体积分型守恒方程的应用}

应用积分型守恒方程可以计算一些简单定常流动过程中流体作用在周围物体上的合力与合力矩.因为定常流动中局部导数项等于零，方程中除质量力外其余各项都是面积分，如果已知控制面上一部分流动参数就可以计算另一部分流动参数.所以应用控制体上守恒方程求解实际问题时，应尽量把已知流动参数的面作为控制面的一部分，例如，通常取已知流体流入和流出的边界作为控制面的一部分，然后应用守恒方程来计算流体作用在其余控制面上的合力与合力矩.

\subsection*{1. 应用举例}

\noindent\textbf{例3.5}\quad 理想不可压缩流体流过水平弯管的定常流动，计算流动过程中作用在弯管上的合力.已知：弯管入口面积$A_1$，出口截面积$A_2$，入口流速$U_1$和压强$p_1$，流体密度$\rho$，弯管转角$\alpha$.

\noindent\textbf{解}\quad 假定弯管进出口截面上速度和压强是均匀分布的.首先取如图3.6所示的坐标系和控制体$ABCD$，它包括：进口截面$(AB)$，管壁$A_w$(包含$BC$,$DA$和上下底面)和出口截面$(CD)$.

要计算流体对管壁的作用力，应当利用动量方程.将理想流体的应力张量$\boldsymbol{T}_n=-p\boldsymbol{n}$代入控制体上的动量方程：
\[
\oint_\Sigma \rho\boldsymbol{U}(\boldsymbol{U}\cdot\boldsymbol{n})\,\mathrm{d}A=\int_D \rho\boldsymbol{f}\,\mathrm{d}V-\oint_\Sigma p\boldsymbol{n}\,\mathrm{d}A
\]
其中动量输运量
\begin{align*}
\oint_\Sigma \rho\boldsymbol{U}(\boldsymbol{U}\cdot\boldsymbol{n})\,\mathrm{d}A&=\iint_{AB}\rho\boldsymbol{U}_1(\boldsymbol{U}_1\cdot\boldsymbol{n}_1)\,\mathrm{d}A\\
&\quad+\iint_{A_w}\rho\boldsymbol{U}(\boldsymbol{U}\cdot\boldsymbol{n})\,\mathrm{d}A+\iint_{CD}\rho\boldsymbol{U}_2(\boldsymbol{U}_2\cdot\boldsymbol{n}_2)\,\mathrm{d}A
\end{align*}
应用连续方程$U_2A_2=U_1A_1$，得$U_2=U_1A_1/A_2$，从而进出口速度向量分别为
\[
\boldsymbol{U}_2=\frac{U_1A_1}{A_2}\boldsymbol{n}_2,\qquad \boldsymbol{U}_1=-U_1\boldsymbol{n}_1
\]
在管壁面$A_w$上$\boldsymbol{U}\cdot\boldsymbol{n}=0$，故控制体上动量的输运量为
\[
\oint_\Sigma \rho\boldsymbol{U}(\boldsymbol{U}\cdot\boldsymbol{n})\,\mathrm{d}A=-\rho U_1U_1A_1+\rho U_2U_2A_2=\rho U_1^{2}A_1\boldsymbol{n}_1+\rho U_2^{2}A_2\boldsymbol{n}_2
\]
动量守恒方程中压强积分为
\[
-\oint_\Sigma p\boldsymbol{n}\,\mathrm{d}A=-\iint_{AB}p_1\boldsymbol{n}\,\mathrm{d}A-\iint_{A_w}p\boldsymbol{n}\,\mathrm{d}A-\iint_{CD}p_2\boldsymbol{n}\,\mathrm{d}A
\]
式中$-\displaystyle\iint_{A_w}p\boldsymbol{n}\,\mathrm{d}A$是管壁作用在流体上的合力，根据作用和反作用原理，它等于流体作用在管壁上力$\boldsymbol{F}$的负值，即
\[
-\iint_{A_w}p\boldsymbol{n}\,\mathrm{d}A=-\boldsymbol{F}
\]
假定进出口截面上压强分布是均匀的，于是压强积分项
\[
-\oint_\Sigma p\boldsymbol{n}\,\mathrm{d}A=-p_1A_1\boldsymbol{n}_1-p_2A_2\boldsymbol{n}_2-\boldsymbol{F}
\]
最后体积力$\displaystyle\int\rho\boldsymbol{g}\,\mathrm{d}V=-\rho gV\boldsymbol{k}$，它是水平弯管内的水重.所有各项代入动量守恒方程后，得弯管上受到的流体作用力
\[
\boldsymbol{F}=-\rho gV\boldsymbol{k}-(p_1+\rho U_1^{2})A_1\boldsymbol{n}_1-(p_2+\rho U_2^{2})A_2\boldsymbol{n}_2
\]
这表明弯管上受到一垂直力，它等于弯管内水的重力，另外有一水平力，它由流体流经弯管时动量变化和流入以及流出截面上的压强合力产生.对于不可压缩理想流体的定常流动，上式中$p_2$可用伯努利方程求出如下：
\[
p_2=p_1+\frac{\rho}{2}(U_1^{2}-U_2^{2})=p_1+\frac{\rho U_1^{2}}{2}\Big[1-\Big(\frac{A_1}{A_2}\Big)^{2}\Big]
\]
代入流体作用力公式就可得作用在管壁的合力.

若$A_1=A_2$，由连续方程和伯努利公式可得：$U_1=U_2=U$和$p_1=p_2$，流体作用力公式可简化为
\[
\boldsymbol{F}=-(p_1A+\rho U^{2}A)\boldsymbol{n}_1-(p_1A+\rho U^{2}A)\boldsymbol{n}_2-\rho gV\boldsymbol{k}
\]
对于图3.3的弯管，$\boldsymbol{n}_1=-\boldsymbol{j}$，$\boldsymbol{n}_2=\boldsymbol{i}\sin\alpha+\boldsymbol{j}\cos\alpha$，代入上式，得
\[
\boldsymbol{F}=-(p_1A+\rho U^{2}A)\boldsymbol{i}\sin\alpha+(p_1A+\rho U^{2}A)(1-\cos\alpha)\boldsymbol{j}-\rho gV\boldsymbol{k}
\]

\noindent\textbf{例3.6}\quad 转角为$90^{\circ}$的弯管，其直径为$1\,$m，通过的水流量为$11500\,\mathrm{m^3/h}$，求水流作用在弯管上的力(不计重力和压强合力).

\noindent\textbf{解}\quad 应用例3.5中公式，$\rho=1000\,\mathrm{kg/m^3}$，$A_1=A_2=A=\dfrac{\pi}{4}\,d^{2}=0.7854\,\mathrm{m^2}$，
$Q=11500/3600=3.194\,(\mathrm{m^3/s})$，$U_1=U_2=U=Q/A=4.067\,$m/s，得
\[
F_x=-12992\,\mathrm{N},\qquad F_y=12992\,\mathrm{N}
\]

\noindent\textbf{例3.7}\quad 有一股水平定常平面理想射流冲击固定平板(图3.7)，已知射流速度为$\boldsymbol{U}$、宽度为$d$、密度为$\rho$、环境压强为$p_a$、平板的法向与射流交角为$\alpha$.求作用在单位厚度平板上的合力及合力作用点(不计质量力).

\noindent\textbf{解}\quad 射流冲击平板后将分成两股，设下游两股流动的宽度分别为$d_1$、$d_2$(见图3.7)，假定入口射流速度分布和下游两股流动在足够远的出口处的速度分布都是均匀的.

首先取直角坐标系如图3.7所示.$y$轴平行于平板壁面，坐标原点为进口射流和平壁面的交点.并取射流表面$BC$,$AF$，进出口截面$AB$,$CD$,$FE$及板面$DE$间流场为控制体(图3.7中的虚线)，在该控制体上应用动量方程，令$\boldsymbol{f}=\boldsymbol{0}$，得
\[
\oint_\Sigma \rho\boldsymbol{U}(\boldsymbol{U}\cdot\boldsymbol{n})\,\mathrm{d}A=-\oint_\Sigma p\boldsymbol{n}\,\mathrm{d}A
\]
动量输运量为
\begin{align*}
\oint_\Sigma \rho\boldsymbol{U}(\boldsymbol{U}\cdot\boldsymbol{n})\,\mathrm{d}A&=\iint_{AB}\rho\boldsymbol{U}(\boldsymbol{U}\cdot\boldsymbol{n})\,\mathrm{d}A+\iint_{BC+FA}\rho\boldsymbol{U}(\boldsymbol{U}\cdot\boldsymbol{n})\,\mathrm{d}A\\
&\quad+\iint_{CD}\rho\boldsymbol{U}_1(\boldsymbol{U}_1\cdot\boldsymbol{n})\,\mathrm{d}A+\iint_{EF}\rho\boldsymbol{U}_2(\boldsymbol{U}_2\cdot\boldsymbol{n})\,\mathrm{d}A
\end{align*}
射流侧面是流面，因而$BC+FA$面上$\boldsymbol{U}\cdot\boldsymbol{n}=0$，在平板面上也有$\boldsymbol{U}\cdot\boldsymbol{n}=0$.平面流中取垂直图面方向为单位厚度，这时动量输运量为
\begin{align*}
\oint_\Sigma \rho\boldsymbol{U}(\boldsymbol{U}\cdot\boldsymbol{n})\,\mathrm{d}A&=-\rho U\,|\boldsymbol{U}|\,d+\rho U_1^{2}d_1\boldsymbol{j}-\rho U_2^{2}d_2\boldsymbol{j}\\
&=-\rho U^{2}d(\boldsymbol{i}\cos\alpha+\boldsymbol{j}\sin\alpha)+\rho U_1^{2}d_1\boldsymbol{j}-\rho U_2^{2}d_2\boldsymbol{j}
\end{align*}
这里$U_1$、$U_2$为未知量，可应用伯努利公式
\[
p_{\mathrm{a}}+\frac{\rho U^{2}}{2}=p_{\mathrm{a}}+\frac{\rho U_1^{2}}{2}=p_{\mathrm{a}}+\frac{\rho U_2^{2}}{2}
\]
于是可求出：$U_1=U_2=U$；另外根据连续性方程：$Ud=U_1d_1+U_2d_2=U(d_1+d_2)$，因而
\[
d=d_1+d_2
\]
再计算压强合力：
\[
-\oint_\Sigma p\boldsymbol{n}\,\mathrm{d}A=-\iint_{AB+BC+CD+EF+FA}p_{\mathrm{a}}\boldsymbol{n}\,\mathrm{d}A-\iint_{DE}p\boldsymbol{n}\,\mathrm{d}A
\]
其中$-\displaystyle\iint_{DE}p\boldsymbol{n}\,\mathrm{d}A$为固壁作用于流体控制体上的合力，射流对固壁的作用力为$\displaystyle\iint_{DE}p\boldsymbol{n}\,\mathrm{d}A$，在固壁的另一侧作用有大气压强，故固壁上总的合力为
\[
\boldsymbol{F}=\iint_{DE}(p-p_{\mathrm{a}})\boldsymbol{n}\,\mathrm{d}A
\]
将压强分布代入上式，得
\begin{align*}
\boldsymbol{F}=\iint_{DE}p\boldsymbol{n}\,\mathrm{d}A-\iint_{DE}p_{\mathrm{a}}\boldsymbol{n}\,\mathrm{d}A&=\oint_\Sigma p\boldsymbol{n}\,\mathrm{d}A-\iint_{AB+BC+CD+EF+FA}p_{\mathrm{a}}\boldsymbol{n}\,\mathrm{d}A-\iint_{DE}p_{\mathrm{a}}\boldsymbol{n}\,\mathrm{d}A\\
&=\oint_\Sigma p\boldsymbol{n}\,\mathrm{d}A-\oint_\Sigma p_{\mathrm{a}}\boldsymbol{n}\,\mathrm{d}A
\end{align*}
最后一个积分是常压强在封闭体表面的合力，它应等于零，所以平板上内外两侧的合力$\boldsymbol{F}=\displaystyle\oint_\Sigma p\boldsymbol{n}\,\mathrm{d}\boldsymbol{A}$.将它和动量输运公式一起代入动量方程后有
\[
-\rho U^{2}d(\boldsymbol{i}\cos\alpha+\boldsymbol{j}\sin\alpha)+\rho U^{2}d_1\boldsymbol{j}-\rho U^{2}d_2\boldsymbol{j}=-\boldsymbol{F}
\]
或
\[
\boldsymbol{F}=(\rho U^{2}d\cos\alpha)\boldsymbol{i}-(\rho U^{2}d_1-\rho U^{2}d_2-\rho U^{2}d\sin\alpha)\boldsymbol{j}
\]
理想流体运动作用在壁面上只有垂直壁面的压强，没有平行壁面的力，故$\boldsymbol{F}$的$y$方向分量等于零，即上式右边第二项等于零，于是应有
\[
\boldsymbol{F}=\rho U^{2}d\cos\alpha\,\boldsymbol{i}\,,\qquad d_1-d_2=d\sin\alpha
\]
联立$d_1+d_2=d$，得到分叉射流的厚度分别为
\[
d_1=d(1+\sin\alpha)/2,\qquad d_2=d(1-\sin\alpha)/2
\]
下面来求合力作用点.这时需要应用动量矩方程
\[
\oint_\Sigma \rho\boldsymbol{r}\times\boldsymbol{U}(\boldsymbol{U}\cdot\boldsymbol{n})\,\mathrm{d}A=-\oint_\Sigma p(\boldsymbol{r}\times\boldsymbol{n})\,\mathrm{d}A
\]
对坐标原点取矩，计算动量矩输运量
\begin{align*}
\oint_\Sigma \rho\boldsymbol{r}\times\boldsymbol{U}(\boldsymbol{U}\cdot\boldsymbol{n})\,\mathrm{d}A&=\iint_{AB}\rho\boldsymbol{r}\times\boldsymbol{U}(\boldsymbol{U}\cdot\boldsymbol{n})\,\mathrm{d}A+\iint_{BC+DE+FA}\rho\boldsymbol{r}\times\boldsymbol{U}(\boldsymbol{U}\cdot\boldsymbol{n})\,\mathrm{d}A\\
&\quad+\iint_{CD}\rho\boldsymbol{r}\times\boldsymbol{U}_1(\boldsymbol{U}_1\cdot\boldsymbol{n})\,\mathrm{d}A+\iint_{EF}\rho\boldsymbol{r}\times\boldsymbol{U}_2(\boldsymbol{U}_2\cdot\boldsymbol{n})\,\mathrm{d}A
\end{align*}
在射流表面和壁面上$\boldsymbol{U}\cdot\boldsymbol{n}=0$，故输运量为零；入口射流的动量矩也为零，因入口速度通过坐标原点，出口两截面上流动的动量矩分别为
\[
\iint_{FE}\rho\boldsymbol{r}\times\boldsymbol{U}_2(\boldsymbol{U}_2\cdot\boldsymbol{n})\,\mathrm{d}A=-\frac{\rho U^{2}d_2^{2}}{2}\boldsymbol{k}\,,\qquad
\iint_{CD}\rho\boldsymbol{r}\times\boldsymbol{U}_1(\boldsymbol{U}_1\cdot\boldsymbol{n})\,\mathrm{d}A=\frac{\rho U^{2}d_1^{2}}{2}\boldsymbol{k}
\]
总的动量矩输运量为
\[
\oint_\Sigma \rho\boldsymbol{r}\times\boldsymbol{U}(\boldsymbol{U}\cdot\boldsymbol{n})\,\mathrm{d}A=-\frac{\rho U^{2}}{2}(d_2^{2}-d_1^{2})\boldsymbol{k}
\]
压强作用力矩可写成
\[
-\oint_\Sigma p\boldsymbol{r}\times\boldsymbol{n}\,\mathrm{d}A=-\oint_\Sigma(p-p_{\mathrm{a}})\boldsymbol{r}\times\boldsymbol{n}\,\mathrm{d}A-\oint_\Sigma p_{\mathrm{a}}\boldsymbol{r}\times\boldsymbol{n}\,\mathrm{d}A
\]
其中最后一积分等于零，右边第一个积分为
\[
-\oint_\Sigma(p-p_{\mathrm{a}})\boldsymbol{r}\times\boldsymbol{n}\,\mathrm{d}A=-\iint_{AB+BC+CD+EF+FA}(p-p_{\mathrm{a}})\boldsymbol{r}\times\boldsymbol{n}\,\mathrm{d}A-\iint_{DE}(p-p_{\mathrm{a}})\boldsymbol{r}\times\boldsymbol{n}\,\mathrm{d}A
\]
上式第一个积分为零，因为那些面上$p=p_{\mathrm{a}}$，最后一积分为固壁作用在流体上的外力矩，流体作用在固壁上外力矩为它的负值.因此，动量矩方程简化为
\begin{align*}
-\frac{\rho U^{2}}{2}(d_2^{2}-d_1^{2})\boldsymbol{k}&=-\iint_{DE}(p-p_{\mathrm{a}})\boldsymbol{r}\times\boldsymbol{n}\,\mathrm{d}A=-\boldsymbol{r}_0\times\boldsymbol{F}\\
&=-(y_0\boldsymbol{j})\times(\rho U^{2}d^{2}\cos\alpha\,\boldsymbol{i})=y_0\rho U^{2}d\cos\alpha\,\boldsymbol{k}
\end{align*}
由上面的等式，平板上合力作用点位于：
\[
y_0=\frac{d_1^{2}-d_2^{2}}{2d\cos\alpha}=\frac{-d\tan\alpha}{2}\,,\qquad x_0=0
\]
\noindent\textbf{例3.8}\quad 气体引射器的参数计算.

气体引射器是利用一股小流量的高速气流带动大流量的低速气流，又称引射泵，如图3.8所示.1-1截面中心的高速气流$A$引射出低速气流$B$，经过平直段混合后到达2-2截面时参数均匀，忽略壁面摩擦.已知介质为空气，$R=287\,\dfrac{\mathrm{N\cdot m}}{\mathrm{kg\cdot K}}$，绝热指数$\gamma=1.4$，$p_1=9\times 10^{4}\,\mathrm{N/m^2}$，$T_{1A}=250\,$K，$A_2=1\,\mathrm{m^2}$，$T_{1B}=280\,$K，$U_{1B}=10\,$m/s，$U_{1A}=200\,$m/s，$A_{1A}=0.15\,\mathrm{m^2}$，$A_{1B}=0.85\,\mathrm{m^2}$，求截面2-2上的气流参数.

\noindent\textbf{解}\quad 选取图3.8虚线控制体，由空气的状态方程可分别求出两股射流的密度和质量流量$\dot{m}_{1A}$和$\dot{m}_{1B}$：
\begin{gather*}
\rho_{1A}=\frac{p_1}{RT_{1A}}=1.25\,\mathrm{kg/m^3}\\[2pt]
\dot{m}_{1A}=\rho_{1A}U_{1A}A_{1A}=37.62\,\mathrm{kg/s}\\[2pt]
\rho_{1B}=\frac{p_1}{RT_{1B}}=1.12\,\mathrm{kg/m^3}\\[2pt]
\dot{m}_{1B}=\rho_{1B}U_{1B}A_{1B}=9.52\,\mathrm{kg/s}
\end{gather*}
由定常流动控制体上的连续方程可知，流入控制面1-1的质量流量$\dot{m}_{1A}+\dot{m}_{1B}$应等于流出控制面2-2的流量$\dot{m}_2$：
\[
\dot{m}_2=\rho_2U_2A_2=\dot{m}_{1A}+\dot{m}_{1B}=47.14\,\mathrm{kg/s}
\]
由定常流控制体动量方程可以导出，控制面1-1上的压强合力和流入动量通量的和应等于控制面2-2上的压强合力与流出动量通量：
\[
(p_1A_{1A}+\dot{m}_{1A}U_{1A})+(p_1A_{1B}+\dot{m}_{1B}U_{1B})=p_2A_2+\dot{m}_2U_2
\]
故
\[
p_2+47.14U_2=97620\,\mathrm{N/m^2}
\]
最后，应用定常流控制体上的能量方程：流入控制面1-1的热焓和动能通量应等于流出控制面2-2的热焓和动能通量；并考虑到$e+\dfrac{p}{\rho}=h=c_pT=\dfrac{\gamma}{\gamma-1}\dfrac{p}{\rho}$，可得
\[
\dot{m}_{1A}\Big(\frac{\gamma}{\gamma-1}\,\frac{p_1}{\rho_{1A}}+\frac{U_{1A}^{2}}{2}\Big)+\dot{m}_{1B}\Big(\frac{\gamma}{\gamma-1}\,\frac{p_1}{\rho_{1B}}+\frac{U_{1B}^{2}}{2}\Big)=\dot{m}_2\Big(\frac{\gamma}{\gamma-1}\,\frac{p_2}{\rho_2}+\frac{U_2^{2}}{2}\Big)
\]
将已知量代入后，得
\[
3.5\,\frac{p_2}{\rho_2}+\frac{U_2^{2}}{2}=273237\,\mathrm{m^2/s^2}
\]
联立动量方程和能量方程的结果，可得
\begin{gather*}
U_2=38.3\,\mathrm{m/s}\,,\qquad p_2=9.58\times 10^{4}\,\mathrm{N/m^2}\\[2pt]
\rho_2=1.23\,\mathrm{kg/m^3}\,,\qquad T_2=p_2/R\rho_2=271\,\mathrm{K}
\end{gather*}
下面我们再举一个在非惯性坐标系中应用控制面方法的例子.

\noindent\textbf{例3.9}\quad 火箭运动方程.

有一火箭，其初始总质量为$M_0$，发射后以$U_0(t)$的速度垂直向上飞行.相对于火箭的排气速度为$U_j$，单位时间的排气质量为$\dot{m}$，排气压力为$p_j$，排气面积为$A_j$，假定$U_j$、$\dot{m}$和$p_j$均为常数.并设飞行时空气阻力为$\boldsymbol{D}=-D\boldsymbol{k}$，试建立火箭运动微分方程.

\noindent\textbf{解}\quad 将坐标固结在火箭上，并取图3.9中虚线所示的控制体.

排气口的质量流量为$\dot{m}$，因此流出控制体的质量流量为
\[
\oint_\tau \rho(\boldsymbol{U}\cdot\boldsymbol{n})\,\mathrm{d}A=\rho_j U_j A_j=\dot{m}
\]
若控制体内的总质量用$M(t)$表示，则控制体内的质量增长率为
\[
\iiint_\tau \frac{\partial\rho}{\partial t}\,\mathrm{d}V=\frac{\partial}{\partial t}\iiint_\tau \rho\,\mathrm{d}V=\frac{\partial M(t)}{\partial t}
\]
将其代入连续方程得
\[
\frac{\partial M(t)}{\partial t}=-\dot{m}
\]
积分可得
\[
M(t)=-\dot{m}t+c
\]
由$t=0$时，$M=M_0$，可得$c=M_0$，所以
\[
M(t)=M_0-\dot{m}t
\]
火箭加速度为$\boldsymbol{a}_0=\dfrac{\mathrm{d}U_0(t)}{\mathrm{d}t}\boldsymbol{k}$，因此非惯性系的控制体动量方程可写为
\[
\iiint_\tau \rho(\boldsymbol{f}-\boldsymbol{a}_0)\,\mathrm{d}V+\oint_A \boldsymbol{T}_n\,\mathrm{d}A-\oint_A \rho(\boldsymbol{U}\cdot\boldsymbol{n})\boldsymbol{U}\,\mathrm{d}A=\iiint_\tau \frac{\partial}{\partial t}(\rho\boldsymbol{U})\,\mathrm{d}V
\]
式中$\boldsymbol{U}$是流体的相对速度，下面我们分别讨论式中各项
\begin{gather*}
\iiint_\tau \rho(\boldsymbol{f}-\boldsymbol{a}_0)\,\mathrm{d}V=-(g+a_0)\iiint_\tau \rho\,\mathrm{d}V\,\boldsymbol{k}=-(g+a_0)M(t)\boldsymbol{k}\\[2pt]
\oint_A \boldsymbol{T}_n\,\mathrm{d}A=\iint_{A-A_j}\boldsymbol{T}_n\,\mathrm{d}A+\iint_{A_j}-\boldsymbol{n}p_j\,\mathrm{d}A=\iint_{A-A_j}\boldsymbol{T}_n\,\mathrm{d}A+\iint_{A_j}p_j\,\mathrm{d}A\,\boldsymbol{k}\\[2pt]
=\iint_{A-A_j}\boldsymbol{T}_n\,\mathrm{d}A+\iint_{A_j}p_{\mathrm{a}}\boldsymbol{k}\,\mathrm{d}A+\iint_{A_j}(p_j-p_{\mathrm{a}})\,\mathrm{d}A\,\boldsymbol{k}=\boldsymbol{D}+(p_j-p_{\mathrm{a}})A_j\boldsymbol{k}
\end{gather*}
式中$\boldsymbol{D}=\displaystyle\iint_{A-A_j}\boldsymbol{T}_n\,\mathrm{d}A+\iint_{A_j}p_{\mathrm{a}}\,\mathrm{d}A\,\boldsymbol{k}$，定义为控制体上所受的总阻力；单位时间流出控制体的动量通量
\[
\oint_A \rho(\boldsymbol{U}\cdot\boldsymbol{n})\boldsymbol{U}\,\mathrm{d}A=-\iint_{A_j}\rho_j U_j^{2}\boldsymbol{k}\,\mathrm{d}A=-\rho_j U_j^{2}A_j\boldsymbol{k}=-\dot{m}U_j\boldsymbol{k}
\]
在动坐标系中，$\dot{m},p_j$为常数，所以火箭喷气推进是定常流动，其控制体内总动量不变，于是有$\dfrac{\partial}{\partial t}\displaystyle\iiint_\tau \rho\boldsymbol{U}\,\mathrm{d}V=0$，将以上各项代回动量方程后，可得
\[
-(g+a_0)M(t)\boldsymbol{k}-D\boldsymbol{k}+(p_j-p_{\mathrm{a}})A_j\boldsymbol{k}+\dot{m}U_j\boldsymbol{k}=\boldsymbol{0}
\]
最后可得，火箭垂直运动的方程：
\[
a_0=\frac{\mathrm{d}U_0(t)}{\mathrm{d}t}=\frac{\dot{m}U_j+(p_j-p_{\mathrm{a}})A_j-D}{m_0-\dot{m}t}-g
\]

\noindent\textbf{例3.10}\quad 不计空气阻力$(D=0)$，且设$p_j=p_{\mathrm{a}}$，求火箭速度的增长规律.

\noindent\textbf{解}\quad 将$D=0$和$p_j=p_{\mathrm{a}}$代入火箭方程，得到最简单的垂直运动的火箭方程：
\[
\frac{\mathrm{d}U_0}{\mathrm{d}t}=\frac{U_j\,\dot{m}}{M_0-\dot{m}t}-g
\]
对上式积分，得火箭速度
\[
U_0(t)-U_0(0)=U_j\ln\frac{M_0}{M_0-\dot{m}t}-gt=U\ln\frac{M_0}{M(t)}-gt
\]
可以看到，火箭的速度随着喷出质量的增加不断加速，喷出质量越多，火箭速度越高；另一方面，重力场的作用使火箭速度降低.再有，喷气速度越大，火箭的加速度就越大.在以上简化模型中，重力加速度$g$视作常数，事实上，当火箭发射很高时，如人造卫星或宇宙飞船的运载火箭，重力加速度不能视作常数，而应当和垂直高度的平方成反比.

\subsection*{2. 叶轮机械的欧拉公式}

工程中常用的压气机、水泵等机械是利用旋转叶轮对流体做功，以提高流体的动能或压强；另外一类气轮机、水轮机等动力机械，则利用高速或高焓气体推动旋转叶轮做功.这两类机械统称叶轮式流体机械.以流动的形式分类，叶轮式机械又可以分成轴流式和离心式.流体以旋转轴方向通过叶轮的机械称作轴流式；流体在垂直于旋转轴的平面内通过叶轮的机械称作离心式.

图3.10(a)是一个离心式叶轮机械的示意图，叶轮以等角速度$\omega$(旋转轴垂直于纸面)旋转，流体在垂直于$\boldsymbol{\omega}$的平面上通过叶轮.叶轮的入口直径为$d_1$，叶轮上有$Z$个叶片，叶轮出口直径为$d_2$.在旋转的叶轮上，流体流入及流出的相对速度$\boldsymbol{w}_1$、$\boldsymbol{w}_2$与叶轮回周速度反向的交角分别称为入口安装角$\beta_1$和出口安装角$\beta_2$，通过叶轮的流体质量流量为$\dot{M}$.下面用流体力学的积分型方程，来计算叶轮向流体提供的功率.

这是一个叶轮式流体机械的基本问题，实际叶轮中流动比较复杂，需要做一些简化以求得简便的工程计算公式.首先假定叶片很薄又很多，叶片很薄意味着可以不计叶片的厚度，叶片很多意味着可以假定任意两叶片间的流动都相同，即流动参数在圆周方向是相同的.此外，还假定流动是平面流，即流动参数在轴向是不变的.

首先要选择恰当的坐标系和控制体，在本例中如果选用固结在地面的坐标系，则流动是非定常的，这对应用流体动力学的控制体方程很不方便.如选用与叶轮一起旋转的坐标，这时在叶片间的流动相对于该坐标系是定常的.所以应当在固结于叶轮上的旋转坐标系中研究叶轮的流动，建立以两叶片的边界和进出口面包围的空间$ABCD$为控制体，如图3.10(b)所示.由于流动是平面的，在轴向取单位长度.

在旋转坐标系中流体通过叶轮的相对速度用$\boldsymbol{w}$表示，在旋转坐标系中质量力还应包括惯性力，由于叶轮是定轴等速旋转，牵连加速度只有向心加速度，故惯性力强度为
\[
-\boldsymbol{a}\equiv-\boldsymbol{\omega}\times(\boldsymbol{\omega}\times\boldsymbol{r})-2\boldsymbol{\omega}\times\boldsymbol{w}=\omega^{2}\boldsymbol{r}+2\boldsymbol{w}\times\boldsymbol{\omega}
\]
其中$\boldsymbol{r}$为平面极坐标向量，右边第一项为离心惯性力强度，第二项为科氏惯性力强度.

控制体上的质量守恒方程，或连续方程为
\[
\oint_\Sigma \rho(\boldsymbol{w}\cdot\boldsymbol{n})\,\mathrm{d}A=0
\]
在叶片表面上$\boldsymbol{w}\cdot\boldsymbol{n}=0$，故连续方程为
\[
-\iint_{A_1}\rho_1\boldsymbol{w}_1\cdot\boldsymbol{n}\,\mathrm{d}A=\iint_{A_2}\rho_2\boldsymbol{w}_2\cdot\boldsymbol{n}\,\mathrm{d}A=\dot{M}/Z
\]
其中$A_1=\pi d_1/Z$，$A_2=\pi d_2/Z$，$Z$为叶轮上的叶片数(注意这里将叶片厚度已经忽略)，在极坐标中上式还可写成
\[
\int_0^{\frac{2\pi}{Z}} \rho\boldsymbol{w}\cdot\boldsymbol{r}\,\mathrm{d}\theta=\frac{\dot{M}}{Z}
\]
为求力矩应用动量矩守恒方程，由于相对运动是定常的，局部导数项为零，动量矩守恒方程应为
\[
\oint_\Sigma \rho\boldsymbol{r}\times\boldsymbol{w}(\boldsymbol{w}\cdot\boldsymbol{n})\,\mathrm{d}A=\int_D \rho\boldsymbol{r}\times\boldsymbol{f}\,\mathrm{d}V+\oint_\Sigma \boldsymbol{r}\times(-p\boldsymbol{n})\,\mathrm{d}A
\]
动量矩输运量为
\begin{align*}
\oint_\Sigma \rho\boldsymbol{r}\times\boldsymbol{w}(\boldsymbol{w}\cdot\boldsymbol{n})\,\mathrm{d}A&=\iint_{A_1}\rho_1\boldsymbol{r}_1\times\boldsymbol{w}_1(\boldsymbol{w}_1\cdot\boldsymbol{n}_1)\,\mathrm{d}A\\
&\quad+\iint_{A}\rho\boldsymbol{r}\times\boldsymbol{w}(\boldsymbol{w}\cdot\boldsymbol{n})\,\mathrm{d}A+\iint_{A_2}\rho_2\boldsymbol{r}_2\times\boldsymbol{w}_2(\boldsymbol{w}_2\cdot\boldsymbol{n}_2)\,\mathrm{d}A
\end{align*}
叶片表面$A$上$\boldsymbol{w}\cdot\boldsymbol{n}=0$；在进出口面上流动是均匀的，因此$A_1$，$A_2$面上的积分式中$\rho_1\boldsymbol{r}_1\times\boldsymbol{w}_1=-\rho_1r_1w_1\cos\beta_1\,\boldsymbol{k}$，$\rho_2\boldsymbol{r}_2\times\boldsymbol{w}_2=-\rho_2r_2w_2\cos\beta_2\,\boldsymbol{k}$，代入积分式后，得
\begin{gather*}
\iint_{A_1}\rho_1\boldsymbol{r}_1\times\boldsymbol{w}_1(\boldsymbol{w}_1\cdot\boldsymbol{n}_1)\,\mathrm{d}A=\boldsymbol{r}_1\times\boldsymbol{w}_1\bigg[\iint_{A_1}\rho_1(\boldsymbol{w}_1\cdot\boldsymbol{n}_1)\,\mathrm{d}A\bigg]\\[4pt]
=-\boldsymbol{r}_1\times\boldsymbol{w}_1\Big(\frac{\dot{M}}{Z}\Big)=\frac{\dot{M}}{Z}\,r_1w_1\cos\beta_1\,\boldsymbol{k}\\[4pt]
\iint_{A_2}\rho_2\boldsymbol{r}_2\times\boldsymbol{w}_2(\boldsymbol{w}_2\cdot\boldsymbol{n}_2)\,\mathrm{d}A=\frac{(\boldsymbol{r}_2\times\boldsymbol{w}_2)\,\dot{M}}{Z}=-\frac{\dot{M}}{Z}\,r_2w_2\cos\beta_2\,\boldsymbol{k}
\end{gather*}
总的动量矩输运量为
\[
\oint_\Sigma \rho\boldsymbol{r}\times\boldsymbol{w}(\boldsymbol{w}\cdot\boldsymbol{n})\,\mathrm{d}A=\frac{\dot{M}}{Z}(r_1w_1\cos\beta_1-r_2w_2\cos\beta_2)\boldsymbol{k}
\]
下面计算质量力矩.
\[
\int_D \rho\boldsymbol{r}\times\boldsymbol{f}\,\mathrm{d}V=\int_D \rho\boldsymbol{r}\times(\omega^{2}\boldsymbol{r}-2\boldsymbol{\omega}\times\boldsymbol{w})\,\mathrm{d}V
\]
第一项$\boldsymbol{r}\times(\omega^{2}\boldsymbol{r})=\boldsymbol{0}$，第二项$-2\boldsymbol{r}\times(\boldsymbol{\omega}\times\boldsymbol{w})=-2[(\boldsymbol{r}\cdot\boldsymbol{w})\boldsymbol{\omega}-(\boldsymbol{r}\cdot\boldsymbol{\omega})\boldsymbol{w}]$，因平面向径$\boldsymbol{r}$与旋转角速度$\boldsymbol{\omega}$垂直，故$\boldsymbol{r}\cdot\boldsymbol{\omega}=0$，因此质量力矩为
\[
\int_D \rho\boldsymbol{r}\times\boldsymbol{f}\,\mathrm{d}V=-2\int_D \rho(\boldsymbol{r}\cdot\boldsymbol{w})\boldsymbol{\omega}\,\mathrm{d}V\,\boldsymbol{k}=-2\boldsymbol{\omega}\boldsymbol{k}\int_D \rho(\boldsymbol{r}\cdot\boldsymbol{w})\,\mathrm{d}V
\]
将体积元$\mathrm{d}V$写成柱坐标式，并在轴向取单位长度，即$\mathrm{d}V=r\,\mathrm{d}r\,\mathrm{d}\theta$，得质量力矩的积分式为
\[
\int_D \rho\boldsymbol{r}\times\boldsymbol{f}\,\mathrm{d}V=-2\int_{r_1}^{r_2}\mathrm{d}r\int \rho r(\boldsymbol{r}\cdot\boldsymbol{w})\omega\,\mathrm{d}\theta\,\boldsymbol{k}=-2\omega\boldsymbol{k}\int_{r_1}^{r_2}r\,\mathrm{d}r\int_0^{\frac{2\pi}{Z}} \rho(\boldsymbol{r}\cdot\boldsymbol{w})\,\mathrm{d}\theta
\]
利用质量连续方程：
\[
\int_0^{\frac{2\pi}{Z}} \rho(\boldsymbol{r}\cdot\boldsymbol{w})\,\mathrm{d}\theta=\int_0^{\frac{2\pi}{Z}} \rho(\boldsymbol{e}_r\cdot\boldsymbol{w})\,r\,\mathrm{d}\theta=\frac{\dot{M}}{Z}
\]
将它代入质量力矩积分式，最后得质量力矩：
\[
\int_D \rho\boldsymbol{r}\times\boldsymbol{f}\,\mathrm{d}V=\frac{-2\omega\,\dot{M}\boldsymbol{k}}{Z}\int_{r_1}^{r_2}r\,\mathrm{d}r=\frac{-\omega\,\dot{M}\boldsymbol{k}}{Z}(r_2^{2}-r_1^{2})
\]
作用在叶片上的表面力矩为
\[
\oint_\Sigma \boldsymbol{r}\times(-p\boldsymbol{n})\,\mathrm{d}A=\iint_{A_1}\boldsymbol{r}_1\times(-p\boldsymbol{n}_1)\,\mathrm{d}A+\iint_{A_2}\boldsymbol{r}_2\times(-p\boldsymbol{n}_2)\,\mathrm{d}A+\iint_{A}\boldsymbol{r}\times(-p\boldsymbol{n})\,\mathrm{d}A
\]
由于在进出口$A_1$，$A_2$上，$\boldsymbol{n}_1=-\boldsymbol{e}_r$，$\boldsymbol{n}_2=\boldsymbol{e}_r$，故$\boldsymbol{r}_1\times\boldsymbol{n}_1=\boldsymbol{0}$，$\boldsymbol{r}_2\times\boldsymbol{n}_2=\boldsymbol{0}$.因而进出口面上压强力矩为零.在叶片表面上积分$\displaystyle\iint \boldsymbol{r}\times(-p\boldsymbol{n})\,\mathrm{d}A$等于叶片作用在流体上的力矩，写作$(L/Z)\boldsymbol{e}_z$，$L$为叶轮对流体作用的总力矩值，即
\[
\oint_\Sigma \boldsymbol{r}\times(-p\boldsymbol{n})\,\mathrm{d}A=\frac{L}{Z}\boldsymbol{k}
\]
将总动量矩输运量、质量力矩和叶轮对流体作用的总力矩值代入动量矩方程，得
\[
\frac{\dot{M}}{Z}(r_1w_1\cos\beta_1-r_2w_2\cos\beta_2)\boldsymbol{k}=\frac{L}{Z}\boldsymbol{k}-\frac{\dot{M}}{Z}(r_2^{2}-r_1^{2})\omega\boldsymbol{k}
\]
简化后，得叶轮对流体作用的总力矩等于：
\begin{equation*}
L=\dot{M}\big[r_1w_1\cos\beta_1-\omega r_1^{2}-(r_2w_2\cos\beta_2-\omega r_2^{2})\big]
\tag{3.23}
\end{equation*}
叶轮对流体所做功率等于作用力矩乘以角速度，即：$P=L\omega$，代入式(3.23)，得
\begin{equation*}
P=\dot{M}\omega\big[r_1w_1\cos\beta_1-\omega r_1^{2}-(r_2w_2\cos\beta_2-\omega r_2^{2})\big]
\tag{3.24}
\end{equation*}
式(3.23)，式(3.24)称作欧拉叶轮机械公式，式中$\beta_1,\beta_2,r_1,r_2$是叶轮的几何参数，$\dot{M}$是质量流量，$\omega$是叶轮转速，这些量给定后，只要算出$w_1$、$w_2$便能得到叶片对流体作用的力矩和功率.

在叶轮机械中，常用速度三角形方法来求叶片的进出口相对速度.在叶轮中流体质点的相对速度是$\boldsymbol{w}$；流体质点相对于固定坐标系的绝对速度用$\boldsymbol{c}$表示；叶轮上任意点的速度是圆周速度$\boldsymbol{\omega}\times\boldsymbol{r}$，也就是旋转坐标系的牵连速度.根据质点运动学原理，质点的绝对速度等于相对速度和牵连速度之和：
\begin{equation*}
\boldsymbol{c}=\boldsymbol{w}+\boldsymbol{\omega}\times\boldsymbol{r}
\tag{3.25}
\end{equation*}
用几何方法表示式(3.25)的速度关系式称为速度三角形，见图3.11.

用$\boldsymbol{u}$表示圆周速度$\boldsymbol{\omega}\times\boldsymbol{r}$，它在圆周的切向，叶片进出口处的切线方向与叶轮圆周速度反向的交角$\beta$定义为叶片安装角.在进出口处，流体绝对速度$\boldsymbol{c}$与叶轮圆周速度正向交角定义为进气角和出气角，并用$\alpha_1$、$\alpha_2$表示.由图3.11所示的速度三角形几何关系，可导出
\[
u_1-w_1\cos\beta_1=r_1\omega-w_1\cos\beta_1=c_1\cos\alpha_1=c_{1u}
\]
同理
\[
u_2-w_2\cos\beta_2=r_2\omega-w_2\cos\beta_2=c_2\cos\alpha_2=c_{2u}
\]
代入式(3.23)、式(3.24)后，欧拉叶轮功率公式可改写为
\begin{gather*}
L=\dot{M}(r_2c_2\cos\alpha_2-r_1c_1\cos\alpha_1)\tag{3.26a}\\[2pt]
P=L\omega=\dot{M}(u_2c_{2u}-u_1c_{1u})\tag{3.26b}
\end{gather*}
$c_{1u}$,$c_{2u}$为进出口流体绝对速度的切向分量，$u_1,u_2$为进出口处的轮周速度.

如果$P>0$，叶轮对流体做功，这种机械称为压缩机，压缩机的条件是
\[
u_2c_{2u}-u_1c_{1u}>0
\]
如果$P<0$，流体对叶轮做功，这种机械称为膨胀机，或涡轮机.这时
\[
u_2c_{2u}-u_1c_{1u}<0
\]

\noindent\textbf{例3.11}\quad 设有风机叶轮的内径$d_1=12.5\,$cm，外径$d_2=30\,$cm，叶片宽度$b=2.5\,$cm，转速$n=1725\,$r/min，体积流量$Q=372\,\mathrm{m^3/h}$.空气沿径向流入叶轮$(\alpha_1=90^{\circ})$，进口压强$p_1=97000\,\mathrm{N/m^2}$，气温$T_1=293\,$K，气流出气角$\beta_2=30^{\circ}$(假设流体是理想不可压缩的)

(1)计算$c_1,u_1,w_1,\beta_1$和$c_2,u_2,w_2,\alpha_2$；

(2)求所需力矩$L$和功率$P$.

\noindent\textbf{解}\quad (1)由已知条件计算进出口的牵连速度
\begin{gather*}
\omega=\frac{2\pi n}{60}=180.6\,\mathrm{rad/s}\\[2pt]
u_1=\omega\cdot\frac{d_1}{2}=11.29\,\mathrm{m/s}\\[2pt]
u_2=\omega\cdot\frac{d_2}{2}=27.1\,\mathrm{m/s}
\end{gather*}
再由状态方程计算流体的密度
\[
\rho=\frac{p_1}{RT_1}=\frac{9.7\times 10^{4}}{287\times 293}=1.155\,(\mathrm{kg/m^3})
\]
质量流量为
\[
\dot{M}=\rho Q=\frac{372\times 1.155}{3600}=0.119\,(\mathrm{kg/s})
\]
下面计算进出口的流动参数与几何参数
\[
c_1=\frac{Q}{A}=\frac{Q}{b\pi d_1}=\frac{372}{3600\times 0.025\times 0.125\pi}=10.53\,(\mathrm{m/s})
\]
因为$\alpha=90^{\circ}$，故$\beta_1=\arctan(c_1/u_1)=43^{\circ}$，
\[
w_1=\frac{u_1}{\cos\beta_1}=15.44\,(\mathrm{m/s})
\]
由出口流量，可得
\begin{gather*}
w_2\sin\beta_2=\frac{Q}{A_2}=\frac{Q}{b\pi d_2}=\frac{372}{3600\times 0.025\times 0.3\pi}=4.386\,(\mathrm{m/s})\\[2pt]
w_2=4.386/\sin\beta_2=8.77\,(\mathrm{m/s})\\[2pt]
c_2=\sqrt{u_2^{2}+w_2^{2}-2u_2w_2\cos\beta_2}=20\,(\mathrm{m/s})
\end{gather*}
因为$c_2/\sin\beta_2=w_2/\sin\alpha_2$，所以$\alpha_2=12.67^{\circ}$.

(2)计算扭矩$L$和功率$P$，因为$c_{1u}=0\,(\alpha_1=90^{\circ})$，所以，
\begin{gather*}
L=\dot{M}(c_{2u}r_2-c_{1u}r_1)=0.119\Big(c_2\cos\alpha_2\times\frac{d_2}{2}\Big)=0.35\,(\mathrm{N\cdot m})\\[2pt]
P=\dot{M}\omega=0.35\times 180.6=63.2\,(\mathrm{N\cdot m/s})
\end{gather*}
以上例子说明了应用流体动力学积分型方程计算流体与周围物体的作用力和力矩的方法.现将其要点归纳如下.

(1)首先选择恰当的坐标系，使得在该坐标系中流动是相对定常的；

(2)选择适当的控制体，使得控制体界面上包括要求的未知量和尽可能多的已知量；

(3)利用沿流管的连续方程和伯努利方程求出流入或流出边界上未知流动参数；

(4)计算出动量输运量、动量矩输运量、质量力的合力或合力矩；

(5)由动量或动量矩方程求出流体和周围物体之间的作用力或力矩.

利用积分型守恒方程可以求出流体作用在物体上的合力与合力矩，如果需要知道速度场细节和物体上的应力分布时，积分型方程就不能满足要求了，而需要建立求解流场和应力场的微分方程.
\section*{3.4  流体动力学的微分型控制方程}\addcontentsline{toc}{section}{3.4 流体动力学的微分型控制方程}

微分型控制方程是流体微团的质量守恒和它的动量、能量守恒方程，它们是控制流动的偏微分方程，故又称流动控制方程或支配方程.

\subsection*{1. 连续性方程}

微团的质量守恒方程称为连续性方程.从积分型方程出发来导出该方程.任取一控制体$D$，在控制体上有以下的质量守恒方程：
\[
\int_D \frac{\partial\rho}{\partial t}\,\mathrm{d}V+\oint_\Sigma \rho(\boldsymbol{U}\cdot\boldsymbol{n})\,\mathrm{d}A=0
\]
利用高斯公式(参见附录)将面积分化为体积分，即
\[
\int \rho\boldsymbol{U}\cdot\boldsymbol{n}\,\mathrm{d}A=\int \nabla\cdot(\rho\boldsymbol{U})\,\mathrm{d}V
\]
代入上式，得
\[
\int_D \bigg[\frac{\partial\rho}{\partial t}+\nabla\cdot(\rho\boldsymbol{U})\bigg]\mathrm{d}V=0
\]
在连续流场中上式对任意控制体都成立，则对任意一微元也应成立，故有
\begin{equation*}
\frac{\partial\rho}{\partial t}+\nabla\cdot(\rho\boldsymbol{U})=0
\tag{3.27a}
\end{equation*}
将$\nabla\cdot(\rho\boldsymbol{U})$展开：$\nabla\cdot(\rho\boldsymbol{U})=\boldsymbol{U}\cdot\nabla\rho+\rho\nabla\cdot\boldsymbol{U}$，式(3.27a)还可写成
\[
\frac{\partial\rho}{\partial t}+\boldsymbol{U}\cdot\nabla\rho+\rho\nabla\cdot\boldsymbol{U}=0
\]
或
\begin{equation*}
\frac{D\rho}{Dt}+\rho\nabla\cdot\boldsymbol{U}=0
\tag{3.27b}
\end{equation*}
两种不同形式的连续方程(3.27a)和(3.27b)分别表述不同的物理概念.式(3.27a)中$\nabla\cdot(\rho\boldsymbol{U})$实际上是微元控制体上单位体积的质量通量，因为有高斯公式：
\[
\nabla\cdot(\rho\boldsymbol{U})=\lim_{D\to 0}\frac{\displaystyle\oint_\Sigma \rho(\boldsymbol{U}\cdot\boldsymbol{n})\,\mathrm{d}A}{D}
\]
该式表示散度$\nabla\cdot(\rho\boldsymbol{U})$是流出微元控制体单位体积的质量通量，因此式(3.27a)可表述为

\textbf{微元控制体密度的局部增长率}$\dfrac{\partial\rho}{\partial t}$\textbf{$+$微元控制体单位体积流出的质量}$\nabla\cdot(\rho\boldsymbol{U})=0$

第二种描述式(3.27b)中$\dfrac{D\rho}{Dt}$是微团密度的增长率，$\nabla\cdot\boldsymbol{U}$表示质点的相对体积膨胀率，因此式(3.27b)可表述为

\textbf{微团密度的相对增长率}$\Big(\dfrac{1}{\rho}\dfrac{D\rho}{Dt}\Big)$\textbf{$+$微团的相对体积膨胀率}$(\nabla\cdot\boldsymbol{U})=0$

不可压缩流体的$\dfrac{D\rho}{Dt}=0$，因此它的连续方程可由式(3.27b)简化为
\begin{equation*}
\nabla\cdot\boldsymbol{U}=0
\tag{3.27c}
\end{equation*}
不可压缩流体既可以是均质的，如纯水；也可以是非均质的，例如含有不同盐分的海水，它的密度可以是时间和坐标的函数，但是任一质点的$\dfrac{D\rho}{Dt}$始终等于零，式(3.27c)既适用于均质不可压缩流体，也适用于非均质的不可压缩流体.

\subsection*{2. 运动方程}

流体微团的动量方程称为流体运动方程.仍由积分型动量方程出发来导出该方程，已有积分型动量方程如下：
\[
\int_D \frac{\partial\rho\boldsymbol{U}}{\partial t}\,\mathrm{d}V+\oint_\Sigma \rho\boldsymbol{U}(\boldsymbol{U}\cdot\boldsymbol{n})\,\mathrm{d}A=\iint_\Sigma \boldsymbol{T}_n\,\mathrm{d}A+\int_D \rho\boldsymbol{f}\,\mathrm{d}V
\]
应用高斯公式，将上式中面积分化为体积分：
\begin{gather*}
\iint_\Sigma \rho\boldsymbol{U}(\boldsymbol{U}\cdot\boldsymbol{n})\,\mathrm{d}A=\int_D \nabla\cdot(\rho\boldsymbol{U}\boldsymbol{U})\,\mathrm{d}V=\int_D \nabla\cdot(\rho U_iU_j\boldsymbol{e}_i\boldsymbol{e}_j)\,\mathrm{d}V\\[2pt]
\oint_\Sigma \boldsymbol{T}_n\,\mathrm{d}A=\oint_\Sigma (T_{ij}\boldsymbol{e}_i\boldsymbol{e}_j)\cdot\boldsymbol{n}\,\mathrm{d}A=\int_D \nabla\cdot(T_{ij}\boldsymbol{e}_i\boldsymbol{e}_j)\,\mathrm{d}V
\end{gather*}
将它们代入积分型动量方程，可得
\[
\int_D \bigg[\frac{\partial\rho\boldsymbol{U}}{\partial t}+\nabla\cdot(\rho U_iU_j\boldsymbol{e}_i\boldsymbol{e}_j)-\nabla\cdot(T_{ij}\boldsymbol{e}_i\boldsymbol{e}_j)-\rho\boldsymbol{f}\bigg]\mathrm{d}V=0
\]
上式对任意控制体都成立，在连续流场中上式对任意微元控制体成立，故有
\begin{equation*}
\frac{\partial(\rho\boldsymbol{U})}{\partial t}+\nabla\cdot(\rho U_iU_j\boldsymbol{e}_i\boldsymbol{e}_j)=\rho\boldsymbol{f}+\nabla\cdot(T_{ij}\boldsymbol{e}_i\boldsymbol{e}_j)
\tag{3.28a}
\end{equation*}
注意到$\nabla\cdot T_{ij}\boldsymbol{e}_i\boldsymbol{e}_j=\lim\limits_{D\to 0}\Big(\displaystyle\oint_\Sigma T_{ij}\boldsymbol{e}_i\boldsymbol{e}_j\cdot\boldsymbol{n}\,\dfrac{\mathrm{d}A}{D}\Big)$表示流体微团上单位体积的表面力合力，所以式(3.28a)可表述为

\textbf{微元控制体单位体积流体上的局部动量增长率与通过微元控制体的单位体积流体的动量输出量之和等于微元控制体上单位体积流体的质量力与表面力之和.}

式(3.28a)可以进一步简化.将左边第一、第二项导数分别展开：
\begin{gather*}
\frac{\partial(\rho\boldsymbol{U})}{\partial t}=\boldsymbol{U}\frac{\partial\rho}{\partial t}+\rho\frac{\partial\boldsymbol{U}}{\partial t}\\[2pt]
\nabla\cdot(\rho U_iU_j\boldsymbol{e}_i\boldsymbol{e}_j)=(U_j\boldsymbol{e}_j)\,\nabla\cdot(\rho U_i\boldsymbol{e}_i)+\rho U_i\boldsymbol{e}_i\cdot\nabla U_j\boldsymbol{e}_j\\
=\boldsymbol{U}\nabla\cdot(\rho\boldsymbol{U})+(\rho\boldsymbol{U})\cdot\nabla\boldsymbol{U}
\end{gather*}
代回式(3.28a)，左边两项之和可进一步简化为
\[
\frac{\partial\rho\boldsymbol{U}}{\partial t}+\nabla\cdot(\rho U_iU_j\boldsymbol{e}_i\boldsymbol{e}_j)=\boldsymbol{U}\frac{\partial\rho}{\partial t}+\rho\frac{\partial\boldsymbol{U}}{\partial t}+\boldsymbol{U}\nabla\cdot\rho\boldsymbol{U}+\rho\boldsymbol{U}\cdot\nabla\boldsymbol{U}=\rho\Big(\frac{\partial\boldsymbol{U}}{\partial t}+\boldsymbol{U}\cdot\nabla\boldsymbol{U}\Big)=\rho\,\frac{D\boldsymbol{U}}{Dt}
\]
导出上式时，利用连续性方程$\dfrac{\partial\rho}{\partial t}+\nabla\cdot(\rho\boldsymbol{U})=0$，右边第一、第三项之和等于零.最后微分型的动量方程，或运动方程又可写作：
\begin{equation*}
\rho\,\frac{D\boldsymbol{U}}{Dt}=\rho\boldsymbol{f}+\nabla\cdot(T_{ij}\boldsymbol{e}_i\boldsymbol{e}_j)
\tag{3.28b}
\end{equation*}
注意到$D\boldsymbol{U}/Dt$是质点加速度，$\nabla\cdot T_{ij}\boldsymbol{e}_i\boldsymbol{e}_j/\rho$是流体微团单位质量上表面力的合力，式(3.28b)可表述为

\textbf{流体微团加速度等于作用在微团上单位质量流体的质量力和表面力合力之和.}

\subsection*{3. 能量方程}

流体微团的能量方程也由积分型能量方程出发来导出，已有控制体上的能量方程如下：
\begin{align*}
\int_D \frac{\partial}{\partial t}\bigg[\rho\Big(e+\frac{U^{2}}{2}\Big)\bigg]\mathrm{d}V&+\oint_\Sigma \bigg[\rho\Big(e+\frac{U^{2}}{2}\Big)\boldsymbol{U}\cdot\boldsymbol{n}\bigg]\mathrm{d}A\\
&=\int_D \rho\boldsymbol{f}\cdot\boldsymbol{U}\,\mathrm{d}V+\oint_\Sigma \boldsymbol{U}\cdot\boldsymbol{T}_n\,\mathrm{d}A+\int_D \rho\dot{q}\,\mathrm{d}V+\oint_\Sigma \lambda\,\frac{\partial T}{\partial n}\,\mathrm{d}A
\end{align*}
用高斯公式将上式中的面积分化为体积分：
\begin{gather*}
\oint_\Sigma \rho\Big(e+\frac{U^{2}}{2}\Big)(\boldsymbol{U}\cdot\boldsymbol{n})\,\mathrm{d}A=\int_D \nabla\cdot\bigg[\rho\Big(e+\frac{U^{2}}{2}\Big)\boldsymbol{U}\bigg]\mathrm{d}V\\
=\int_D \rho\boldsymbol{U}\cdot\nabla\Big(e+\frac{U^{2}}{2}\Big)\mathrm{d}V+\int_D \Big(e+\frac{U^{2}}{2}\Big)\nabla\cdot\rho\boldsymbol{U}\,\mathrm{d}V\\[2pt]
\oint_\Sigma \boldsymbol{T}_n\cdot\boldsymbol{U}\,\mathrm{d}A=\oint_\Sigma \boldsymbol{U}\cdot(T_{ij}\boldsymbol{e}_i\boldsymbol{e}_j)\cdot\boldsymbol{n}\,\mathrm{d}A=\int_D \nabla\cdot(\boldsymbol{U}\cdot T_{ij}\boldsymbol{e}_i\boldsymbol{e}_j)\,\mathrm{d}V\\[2pt]
\oint_\Sigma \lambda\,\frac{\partial T}{\partial n}\,\mathrm{d}A=\oint_\Sigma \lambda(\nabla T)\cdot\boldsymbol{n}\,\mathrm{d}A=\int_D \lambda\nabla\cdot\nabla T\,\mathrm{d}V=\int_D \lambda\Delta T\,\mathrm{d}V
\end{gather*}
将以上两项代入能量守恒方程得
\begin{align*}
\int_D \bigg\{\frac{\partial}{\partial t}\bigg[\rho\Big(e+\frac{U^{2}}{2}\Big)\bigg]+\rho\boldsymbol{U}\cdot\nabla\Big(e+\frac{U^{2}}{2}\Big)+\Big(e+\frac{U^{2}}{2}\Big)\nabla\cdot(\rho\boldsymbol{U})\\
-\rho\boldsymbol{f}\cdot\boldsymbol{U}-\nabla\cdot(\boldsymbol{U}\cdot T_{ij}\boldsymbol{e}_i\boldsymbol{e}_j)-\rho\dot{q}-\lambda\Delta T\bigg\}\mathrm{d}V=0
\end{align*}
上式对任意控制体都应成立，因而在连续流场内应处处满足，所以括号内七项和应为零，即
\[
\frac{\partial}{\partial t}\bigg[\rho\Big(e+\frac{U^{2}}{2}\Big)\bigg]+\rho\boldsymbol{U}\cdot\nabla\Big(e+\frac{U^{2}}{2}\Big)+\Big(e+\frac{U^{2}}{2}\Big)\nabla\cdot(\rho\boldsymbol{U})-\rho\boldsymbol{f}\cdot\boldsymbol{U}-\nabla\cdot(\boldsymbol{U}\cdot T_{ij}\boldsymbol{e}_i\boldsymbol{e}_j)-\rho\dot{q}-\lambda\Delta T=0
\]
展开上式第一项局部导数：
\[
\frac{\partial}{\partial t}\bigg[\rho\Big(e+\frac{U^{2}}{2}\Big)\bigg]=\Big(e+\frac{U^{2}}{2}\Big)\frac{\partial\rho}{\partial t}+\rho\,\frac{\partial\big(e+\frac{U^{2}}{2}\big)}{\partial t}
\]
注意到连续方程$\dfrac{\partial\rho}{\partial t}+\nabla\cdot(\rho\boldsymbol{U})=0$，局部导数展开式的第一项与能量方程中第三项可以抵消，最后流体运动的微分型能量方程应为
\begin{equation*}
\frac{\partial}{\partial t}\Big(e+\frac{U^{2}}{2}\Big)+\boldsymbol{U}\cdot\nabla\Big(e+\frac{U^{2}}{2}\Big)=\boldsymbol{f}\cdot\boldsymbol{U}+\frac{1}{\rho}\,\nabla\cdot(\boldsymbol{U}\cdot T_{ij}\boldsymbol{e}_i\boldsymbol{e}_j)+\dot{q}+\frac{\lambda}{\rho}\,\Delta T
\tag{3.29}
\end{equation*}
能量方程(3.29)可表述为

\textbf{微团单位质量流体的能量增长率$=$质量力功率$+$单位质量流体的表面力功率$+$微团内单位质量流体的生成热和微团单位质量流体上由热传导输入的热量之和.}

\subsection*{4. 微分型方程组的封闭性讨论}

前面已经导出了微分型流体运动控制方程(3.27)，(3.28)，(3.29)，把它们归纳如下.

连续性方程
\[
\frac{\partial\rho}{\partial t}+\nabla\cdot(\rho\boldsymbol{U})=0
\]
运动方程
\[
\frac{\partial\boldsymbol{U}}{\partial t}+\boldsymbol{U}\cdot\nabla\boldsymbol{U}=\boldsymbol{f}+\frac{1}{\rho}\,\nabla\cdot(T_{ij}\boldsymbol{e}_i\boldsymbol{e}_j)
\]
能量方程
\[
\frac{\partial\big(e+\frac{U^{2}}{2}\big)}{\partial t}+\boldsymbol{U}\cdot\nabla\Big(e+\frac{U^{2}}{2}\Big)=\boldsymbol{f}\cdot\boldsymbol{U}+\frac{1}{\rho}\,\nabla\cdot(\boldsymbol{U}\cdot T_{ij}\boldsymbol{e}_i\boldsymbol{e}_j)+\dot{q}+\frac{\lambda}{\rho}\,\Delta T
\]
给定适当的初始条件和边界条件后，求解这组微分方程就可以计算出速度场、应力场等.但是方程可解的必要条件是方程数和未知量数必须相等，即方程的封闭性.现在有五个数量方程(动量方程是向量方程，它包含3个数量方程)，而未知数共12个：$\rho,e,T,\boldsymbol{U}$(3个)，$T_{ij}$(6个)，所以方程中未知量的数目大于方程数(缺少7个数量方程)，因而方程组(3.27)，(3.28)，(3.29)是不封闭的.方程组的不封闭性表明需要补充独立方程，这些方程是有别于守恒定理的其他物理规律，例如：热力学的定律、物性方程等.

(1)热力学状态方程

流体微团是宏观热力学平衡体，因此方程中出现的热力学状态参数只有两个独立变量，对于气体，可将内能写成$T$,$\rho$的函数：
\begin{equation*}
e=e(T,\rho)
\tag{3.30}
\end{equation*}
压强、密度和温度之间有状态方程：
\begin{equation*}
p=p(\rho,T)
\tag{3.31}
\end{equation*}
例如，完全气体的状态方程为
\begin{equation*}
p=R\rho T
\tag{3.32a}
\end{equation*}
对于均质不可压缩液体，密度是给定常量，它的状态方程可写作
\begin{equation*}
\rho=\mathrm{const.}
\tag{3.32b}
\end{equation*}
因此介质的状态方程可以补充一个方程.关于热力学状态方程的具体形式和应用，将在后续章节中详细介绍.

(2)本构方程

流体微团的应力状态和微团变形运动状态间的物性关系式称为流体的本构方程.一般情况下，微团应力状态是微团变形运动的泛函：
\begin{equation*}
T_{ij}=T_{ij}(\nabla\boldsymbol{U},\cdots,t)
\tag{3.33}
\end{equation*}
本构方程是张量方程，它有六个代数方程.在求解运动方程组前必须已知本构方程，从而控制方程得以封闭.一般流体的本构方程由统计力学方法或理性力学方法与实验相结合导出，它独立于流体运动方程.在第8章讲授粘性流体力学时将较为详细地讨论本构方程的一般性质及具体形式.补充了热力学状态方程(1个)和本构方程(6个)后，流体力学方程组就封闭了.

最简单的流体，即第1章介绍的理想流体，其应力张量是各向同性的并且和微团的变形运动无关：
\begin{equation*}
T_{ij}=-p\,\delta_{ij}
\tag{3.34}
\end{equation*}
用理想流体的本构方程(3.34)代入运动方程后，未知量减少为7个：$\boldsymbol{U}$(3个)，$p,\rho,e,T$；而理想流体运动的控制方程组恰好有7个方程：1个连续性方程、3个运动方程、1个能量方程、1个热力学状态方程(3.30)和1个应力状态方程(3.34).因此理想流体动力学方程组是封闭的.

\subsection*{5. 边界条件和初始条件}

流体动力学方程组是非线性偏微分方程，需要给定适当的初始和边界条件才能有确定的解，下面讨论常用的几类边界条件和初始条件.

(1)边界条件

边界条件通常可以分为以下几类：

\textcircled{1} 静止无界流场中无穷远条件.例如：航天器在静止大气中飞行，大气的环境参数是已知的，因此远离飞行器的无穷远处流场条件为
\begin{equation*}
|\boldsymbol{x}|\to\infty:\ \boldsymbol{U}=\boldsymbol{0}\,,\quad T_{ij}=-p_\infty\,\delta_{ij}\,,\quad \rho=\rho_\infty\,,\quad T=T_\infty
\tag{3.35}
\end{equation*}
式中$p_\infty$,$\rho_\infty$,$T_\infty$分别是环境的压强、密度和绝对温度.

\textcircled{2} 固壁

在流体中运动的任意固体壁面上，例如在航天器表面上，根据第1章论述的热力学平衡条件，流体既无滑移，又无温差，因此在固壁$\Sigma_b$上流体速度和温度应等于当地固壁的速度和温度：
\begin{equation*}
\Sigma_b:\boldsymbol{U}_f=\boldsymbol{U}_b\,,\quad T_f=T_b
\tag{3.36}
\end{equation*}
式中下标$f$表示流体参数，下标$b$表示固壁的相应物理参数.

\textcircled{3} 互不掺混流体的界面

液体中的气泡、空气中的液滴和海洋表面等都存在两种不同流体的界面.在忽略表面张力情况下，热力学平衡的界面条件为
\begin{equation*}
\Sigma_f:\boldsymbol{U}_{+}=\boldsymbol{U}_{-}\,,\quad T_{+}=T_{-}\,,\quad (T_{ij})_{+}=(T_{ij})_{-}
\tag{3.37}
\end{equation*}
式中$\Sigma_f$表示互不掺混液体界面，下标“$+,-$”表示界面的两侧.流体界面上的边界条件中，除了速度和温度连续外，还有应力张量连续.液体界面和固壁不同，通常液体界面随着流体运动而变化，它的几何方程是未知量，需要将它和运动方程、能量方程联立求解，在水波问题中将讨论流体界面和控制方程联立求解方法.

当表面张力不可忽略时，界面上的剪应力仍然连续，而法向应力差应和表面张力平衡(参见第1章)，这时流体界面的应力边界条件的一般表达式为
\begin{equation*}
\begin{cases}
(\boldsymbol{T}_n\cdot\boldsymbol{n})_{+}-(\boldsymbol{T}_n\cdot\boldsymbol{n})_{-}=\gamma\Big(\dfrac{1}{R_1}+\dfrac{1}{R_2}\Big)\\[8pt]
(\boldsymbol{T}_n\cdot\boldsymbol{s})_{+}=(\boldsymbol{T}_n\cdot\boldsymbol{s})_{-}\,,\quad (\boldsymbol{T}_n\cdot\boldsymbol{t})_{+}=(\boldsymbol{T}_n\cdot\boldsymbol{t})_{-}
\end{cases}
\tag{3.38}
\end{equation*}
式中$\boldsymbol{n}$为流体界面的法向量，$\boldsymbol{s},\boldsymbol{t}$是流体界面的切向量.

(2)初始条件

初始条件是流场的初始状态，一般情况下，在初始时刻$t=t_0$，应给出速度场和热力学状态的分布：
\begin{equation*}
t=0:\ \boldsymbol{U}=\boldsymbol{U}(\boldsymbol{x})\,,\quad p=p(\boldsymbol{x})\,,\quad \rho=\rho(\boldsymbol{x})
\tag{3.39}
\end{equation*}
对于定常流动，流场和时间无关，因此不需要提供初始条件，由定常流动的控制方程(所有局部导数$\partial(\ )/\partial t=0$)和边界条件就可以解出定常流场.

流体运动的基本方程只有一组，但是流体运动却千姿百态，有时奇妙无比，这完全是初始状态和边界条件千差万别的缘故.对于具体问题，给定正确的初始和边界条件是非常重要的，后续章节中，将结合具体问题给出这些条件的数学公式.

\section*{3.5  流体静力学}\addcontentsline{toc}{section}{3.5 流体静力学}

本章最后把静止流场作为微分型流体控制方程的一个应用特例来讨论.根据流体的易流性，静止流体中不存在剪应力，因此它的应力状态是各向同性的：$T_{ij}=-p\,\delta_{ij}$，就是说静止流场中一点的应力状态只有压强.

静止流场中速度处处为零$(\boldsymbol{U}=\boldsymbol{0})$，因此流体静力学要解决的问题是在已知边界和外力场下求静止流场的压强分布，进一步可以求得静止流场中物体上的合力与合力矩.例如静止大气中压强随海拔高度的变化、水利建筑物上的静止压强和合力等.

\subsection*{1. 流体静力学基本方程及积分}

将$\boldsymbol{U}=\boldsymbol{0}$和$T_{ij}=-p\,\delta_{ij}$代入流体动力学基本方程组，此时连续方程(3.27)自然满足，能量方程(3.29)简化为静止介质中的导热方程，运动方程(3.28)简化为流体静力学方程：
\begin{equation*}
-\nabla p+\rho\boldsymbol{f}=\boldsymbol{0}
\tag{3.40a}
\end{equation*}
或写成分量形式：
\begin{equation*}
\frac{\partial p}{\partial x_i}=\rho f_i
\tag{3.40b}
\end{equation*}
如已知外力场$\boldsymbol{f}$和密度场，直接积分式(3.40)就可以求出静止流场的压强分布：
\[
\int_{x_0}^{x}\nabla p\cdot\mathrm{d}\boldsymbol{x}=\int_{x_0}^{x}\frac{\partial p}{\partial x_i}\,\mathrm{d}x_i=\int_{x_0}^{x}\rho f_i\,\mathrm{d}x_i
\]
由于$\dfrac{\partial p}{\partial x_i}\mathrm{d}x_i$是全微分，得
\begin{equation*}
p(\boldsymbol{x})-p(\boldsymbol{x}_0)=\int_{x_0}^{x}\rho f_i\,\mathrm{d}x_i
\tag{3.41}
\end{equation*}
式(3.41)是静止流场中计算压强分布的公式，它的形式相当简单，但是应当指出静力学方程求解的两个定解条件：

(1)右端积分$\displaystyle\int_{x_0}^{x}\rho f_i\,\mathrm{d}x_i$必须和积分路径无关，否则式(3.41)压强是多值的，这在物理上是不可能的，就是说：

\textbf{体积力强度$\rho\boldsymbol{f}$的线积分与积分路径无关是流体能否静止的必要条件.}

(2)外力场满足全微分条件时，流体静力学问题的定解条件是：必须给定一点压强$p(\boldsymbol{x}_0)$.

\subsection*{2. 静止流场中的质量力条件及静止流场的性质}

前面论述了静力学方程的定解条件，下面进一步导出静止流场中的质量力必须满足的条件及静止流场的若干性质.

(1)静止流场的质量力条件

根据斯托克斯定理(见附录)，线积分$\displaystyle\int\rho f_i\,\mathrm{d}x_i=\int\rho\boldsymbol{f}\cdot\mathrm{d}\boldsymbol{x}$与积分路径无关的充要条件是
\[
\nabla\times(\rho\boldsymbol{f})=\boldsymbol{0}
\]
应用向量算符公式：$\nabla\times(\rho\boldsymbol{f})=\nabla\rho\times\boldsymbol{f}+\rho\nabla\times\boldsymbol{f}$，上式可进一步写作
\begin{equation*}
-\boldsymbol{f}\times\nabla\rho+\rho\nabla\times\boldsymbol{f}=\boldsymbol{0}
\tag{3.42}
\end{equation*}
式(3.42)是静止流场中质量力必须满足的必要条件，下面分两种情况讨论.

\textcircled{1} 正压流体

\noindent\textbf{定义3.5}\quad 如果流体质点的密度只是当地压强的单值函数，则称这种流体是正压流体.

例如：均质绝热气体是正压流体，这时压强和密度有以下的单值关系：
\begin{equation*}
\frac{p}{\rho^{\gamma}}=\mathrm{const.}
\tag{3.43}
\end{equation*}
不可压缩均质流体的全场密度等于常数，即$\rho=\mathrm{const.}$，它也具有正压性.

一般正压流体关系式为
\begin{equation*}
\rho=\rho(p)
\tag{3.44}
\end{equation*}
式(3.44)又称正压方程.对于正压流体有以下定理.

\noindent\textbf{定理3.2}\quad 正压流体静止的必要条件是质量力有势.

\noindent\textbf{证明}\quad 由正压方程可以导出$\nabla\rho=\dfrac{\mathrm{d}\rho}{\mathrm{d}p}\nabla p$，把它代入式(3.42)，得
\[
-\frac{\mathrm{d}\rho}{\mathrm{d}p}\,\nabla p\times\boldsymbol{f}+\rho\nabla\times\boldsymbol{f}=\boldsymbol{0}
\]
由于静止时$\boldsymbol{f}=\nabla p/\rho$，上式第一项等于零.因而正压流体静止的必要条件是
\begin{equation*}
\nabla\times\boldsymbol{f}=\boldsymbol{0}
\tag{3.45}
\end{equation*}
外力场的旋度等于零，则外力场必有势，于是证毕.

全场密度为常数的均匀不可压缩流体是最简单的正压流体，它的静止必要条件同上，也是式(3.45).流体静止的必要条件告诉我们，只有在势力场中正压流体或不可压缩流体才可能有静止状态.地球重力场、等速旋转系统的离心力场和静电力场等都是有势力场，在这种力场中正压流体可能静止.

\textcircled{2} 斜压流体

不满足式(3.44)的流体称为斜压流体.

\noindent\textbf{定理3.3}\quad 斜压流体静止的必要条件是
\begin{equation*}
\boldsymbol{f}\cdot(\nabla\times\boldsymbol{f})=0
\tag{3.46}
\end{equation*}
\noindent\textbf{证明}\quad 流体静止必要条件是式(3.42)，$\nabla\rho\times\boldsymbol{f}+\rho\nabla\times\boldsymbol{f}=\boldsymbol{0}$，将该式做$\boldsymbol{f}$的点积，因为$\boldsymbol{f}$垂直于$\nabla\rho\times\boldsymbol{f}$，故$\boldsymbol{f}\cdot(\nabla\rho\times\boldsymbol{f})=0$，于是得斜压流体静止的必要条件：
\[
\boldsymbol{f}\cdot(\nabla\times\boldsymbol{f})=0
\]
证毕.

(2)静止流场的主要性质

\textcircled{1} 有势力场作用下静止流场中等压面、等密面和等势面三者重合.

首先由平衡方程$\nabla p=\rho\boldsymbol{f}$及质量力的公式$\boldsymbol{f}=-\nabla\Pi$，可得有势力场中的平衡方程：
\begin{equation*}
\nabla p=-\rho\nabla\Pi
\tag{3.47}
\end{equation*}
上式两边作$\nabla\Pi$的矢积，由于等式右边$-\rho\nabla\Pi\times\nabla\Pi=\boldsymbol{0}$，故有
\begin{equation*}
\nabla p\times\nabla\Pi=\boldsymbol{0}
\tag{3.48}
\end{equation*}
由向量分析可知，物理量的梯度向量，如$\nabla p$，平行于物理量等值面的法线方向，因此式(3.47)或式(3.48)表明：\textbf{静止流体中等压面和等势面的法向量处处平行}，要满足这一条件，只有当这两个等值面完全重合时才有可能.事实上，只有正压流体才能在势力场中静止.对于正压流体，密度是压强的单值函数$\rho=\rho(p)$，因此可以定义\textbf{压强函数}$\mathscr{P}$：
\begin{equation*}
\mathscr{P}=\int\frac{\mathrm{d}p}{\rho}=\mathscr{P}^{*}(\rho)\quad\text{或}\quad \mathscr{P}=\int\frac{\mathrm{d}p}{\rho}=\mathscr{P}(p)
\tag{3.49}
\end{equation*}
上式还可以写成梯度式：
\[
\nabla\mathscr{P}=\frac{\nabla p}{\rho}
\]
将压强函数$\mathscr{P}$代入静力学方程，得
\[
\nabla\mathscr{P}+\nabla\Pi=0
\]
积分上式，得
\begin{equation*}
\mathscr{P}(p)+\Pi=C
\tag{3.50}
\end{equation*}
$C$是任意常数.上式清楚地说明等压面和等势面重合.

再证明等密面和等势面重合.对于正压流体，压强是密度的单值函数，等压面也就是等密面，因此等密面必和等势面重合.

\textcircled{2} 两种互不掺混的均质流体的静止交界面是等势面.

性质\textcircled{2}可作为性质\textcircled{1}的推论.用反证法证明这一性质，假定具有不同密度的两种液体的交界面为$\Sigma_2$，但是$\Sigma_2$与等势面$\Sigma_1$不重合，它们有交线$A$(图3.12).在介质$\rho_2$中沿等势面$\Sigma_1$由上而下通过$A$线时，在等势面上密度由$\rho_2$变为$\rho_1$，这和在静止介质$\rho_2$中等密面和等势面相重的定理矛盾，因此交界面必和等势面相重.

\subsection*{3. 重力场中静止流体的压强分布}

(1)不可压缩均质流体在重力场中静平衡时其平衡方程为
\begin{gather*}
\nabla p=-\rho\nabla\Pi\tag{3.51}\\
\rho=\mathrm{const.}\tag{3.52}\\
\Pi=gz\tag{3.53}
\end{gather*}
积分静平衡方程(3.51)得
\begin{equation*}
p=\mathrm{const.}-\rho g z
\tag{3.54}
\end{equation*}
式中的积分常数在流体的连通域中不变，$z$坐标与重力加速度方向相反，以向上为正.根据静平衡方程的定解条件，需要给定流场的一点压强，才能确定静止流场的压强分布.通常取液体和大气交界面上的压强作为定解条件.设$z=z_0$为大气与液体的交界面，该处压强为大气压，则常数const.$=p_{\mathrm{a}}+\rho g z_0$，即压强分布为
\begin{equation*}
p=p_{\mathrm{a}}+\rho g\,(z_0-z)
\tag{3.55}
\end{equation*}
重力场中静止液体压强公式适用于液体连通域中任意点，而与盛液体容器的外形以及液体中的物体无关.式(3.55)是重力场液体压强测量和计算作用在液体中固壁上合力或合力矩的基本公式.

(2)倾斜式微压计的压强测量公式

倾斜式微压计是一种常用的气体压强测量仪器(见图3.13)，它由一较大的盛液容器和一倾斜的细管连通组成.容器与大气相通，容器液面上压强为大气压强$p_{\mathrm{a}}$，细管上端与被测压强连接.下面导出倾斜式微压计测压公式.以盛液容器液面为参考坐标面$z_0=0$，应用式(3.55)得
\[
p=p_{\mathrm{a}}-\rho g z
\]
另一方面，$z=l\sin\alpha$，$l$为倾斜管内的液柱升高的长度，$\alpha$为倾斜管与水平面交角，这样在倾斜管端的压强为
\[
p-p_{\mathrm{a}}=-\rho g l\sin\alpha
\]
读出细管中液柱长度，就可以测出当地压强和大气压强之差.这种压强计通常用于测量低于大气压的气体压强，由于倾斜长度$l$大于垂直高度$\Delta Z$，因此可以提高测量的分辨度.例如当$\sin\alpha=0.1$，读数$l=10\,$mm时实际压强与大气压强之差为液柱高$1\,$mm.因此称这种测压计为倾斜式微压计.

实际使用微压计的压差公式需要作修正.通常测量值写作：
\[
p-p_{\mathrm{a}}=-\xi\rho g l\sin\alpha
\]
式中$\xi$称为微压计修正系数.修正系数在仪器使用前用标准压强进行标定，需要修正的原因之一是当微压计读数$l>0$时，由于液体的总体积不变，当倾斜管中液面上升时，盛液容器中液位下降；也就是说，盛液容器中液位(式(3.55))的参考液面$z_0$低于倾斜管上标注的零点.因此倾斜管上标出的液柱长度小于实际压强差应有的液位差(见图3.13).如果盛液容器的横截面远远大于倾斜管的截面，那么这种修正很小.这种修正值的计算并不困难，请读者完成习题3.24的计算.修正的其他原因还有倾斜管中液体的毛细现象等，在微压计使用前可以用标准压差计来标定修正系数$\xi$.

(3)两种不掺混均质液体在重力场中静止时的压强分布

每一均质液体的连通域中压强分布为前面导出的式(3.54)：
\begin{equation*}
p_i=C_i-\rho_i g z
\tag{3.56a}
\end{equation*}
$i$表示不同介质的连通域中的流体参数.设$z_0$为两种液体交界面的液位高，$p_0$为该处的压强，则Ⅰ、Ⅱ两种液体中的压强分布分别为
\begin{gather*}
p_{\mathrm{I}}=p_0-\rho_{\mathrm{I}}\,g(z-z_0)\tag{3.56b}\\[2pt]
p_{\mathrm{II}}=p_0-\rho_{\mathrm{II}}\,g(z-z_0)\tag{3.56c}
\end{gather*}
常识告诉我们处于高液位的液体密度应当低于它下面的液体密度，否则两种液体的静止是不稳定的，即
\begin{equation*}
\rho_{\mathrm{II}}\Big|_{(z<z_0)}>\rho_{\mathrm{I}}\Big|_{(z>z_0)}
\tag{3.57}
\end{equation*}
应用式(3.56a)和式(3.56b)可导出水银压差计测量公式.有水封的水银柱U形管压差计或多管压强计常用来测量水流中的压差或气体中较大的压差.因为水银在空气中会蒸发，它对人体是有害的，因此在水银上方灌以清水.图3.14为一水银U形管的简图.

U形管的两臂所在平面是铅直的，水银上面为水，左侧上端压强为$p_1$，右侧上端压强为$p_2$，水银液位差为$h$.以左端水银液面作为重力势的零点，根据连通域中等压面与等势面重合的性质，在U形管左臂$z=0$面上的压强$p=p_1+\rho_{\text{水}}\,gH$，应等于右臂在$z=0$面上的压强$p=p_2+\rho_{\text{水}}\,g(H-h)+\rho_{\text{汞}}\,gh$，因而有
\begin{equation*}
p_1-p_2=(\rho_{\text{汞}}-\rho_{\text{水}})gh
\tag{3.58}
\end{equation*}
已知$\dfrac{\rho_{\text{汞}}}{\rho_{\text{水}}}=13.6$，代入上式后得
\[
p_1-p_2=12.6\rho_{\text{水}}\,gh
\]

(4)重力场中的静止大气

大气层的压强、密度和温度随高度的变化是航空、气象的重要环境参数.最简单的大气模型假定大气是静止的，并假定重力场是平行力场，即忽略地球的曲率.根据以上假定，大气平衡方程为
\begin{equation*}
\frac{\partial p}{\partial z}=-\rho g\,,\qquad \frac{\partial p}{\partial x}=\frac{\partial p}{\partial y}=0
\tag{3.59}
\end{equation*}
就是说，压强只是高度$z$的函数.由于大气密度与压强有关，式(3.59)自身是不封闭的，要得到式(3.59)的解，必须补充压强与密度的关系式.

\textcircled{1} 绝热大气假定

假定大气层是绝热的，而且大气层中熵不变，这种大气模型称为绝热大气，这时有
\begin{equation*}
\frac{p}{\rho^{\gamma}}=\frac{p_0}{\rho_0^{\gamma}}
\tag{3.60}
\end{equation*}
式中$\gamma$为绝热指数，下标0表示海平面$z=0$上的大气参数.空气接近于双原子分子完全气体，这时$\gamma=1.4$.将式(3.60)代入式(3.59)，并对$z$积分后，可得压强分布：
\begin{equation*}
\frac{\gamma}{\gamma-1}\,\frac{p_0}{\rho_0^{\gamma}}=\frac{\gamma}{\gamma-1}\,\frac{p}{\rho}-gz
\tag{3.61}
\end{equation*}

\textcircled{2} 等温大气假定

假定大气层是等温的，并称之为等温大气，其状态方程为
\begin{equation*}
\frac{p}{\rho}=\frac{p_0}{\rho_0}=RT_0
\tag{3.62}
\end{equation*}
$R$为空气的气体常数，等于$287\,\mathrm{N\cdot m/(kg\cdot K)}$，$T_0$为常数，将式(3.62)代入式(3.59)积分后可得
\begin{equation*}
p=p_0\exp\frac{g(z_0-z)}{RT_0}
\tag{3.63}
\end{equation*}
下标“0”表示参考面$z=z_0$上的大气参数.

\textcircled{3} 国际标准大气

实际大气层中压强和温度关系既非完全绝热也非完全等温.由于近地面大气对流的缘故，大气层温度从地面开始随高度上升而下降，在高度约$11\,$km以上温度几乎不变，称为同温层.国际宇航学会根据以上情况制定了一种国际标准大气，国际标准大气规定的大气层温度分布为
\begin{equation*}
z\leqslant 11000\,\mathrm{m}:\ T=(288-0.0065z)\,\mathrm{K};\qquad z>11000\,\mathrm{m}:\ T=216.5\,\mathrm{K}
\tag{3.64}
\end{equation*}
并称$z\leqslant 11000\,$m为对流层，$z>11000\,$m为同温层.标准大气的压强分布可由给定的温度分布求出.

(i) 对流层中的压强分布

由基本方程：$\dfrac{\mathrm{d}p}{\mathrm{d}z}=-\rho g$，将状态方程：$\rho=\dfrac{p}{RT}=\dfrac{p}{R(288-0.0065z)}$代入，得
\[
\frac{\mathrm{d}p}{p}=\frac{-g\,\mathrm{d}z}{R(288-0.0065z)}
\]
积分后压强分布为
\[
\frac{p}{p_{\mathrm{a}}}=\Big(1-\frac{0.0065z}{288}\Big)^{\frac{g}{0.0065R}}
\]
$p_{\mathrm{a}}$为海平面$z=0$上的大气压强，将物理常数$g$、$R$的数值代入后得对流层中的压强分布：
\begin{equation*}
p=p_{\mathrm{a}}\Big(1-\frac{z}{44300}\Big)^{5.256}\,,\qquad z\leqslant 11000
\tag{3.65}
\end{equation*}
对流层上边界$z=11000\,$m处压强为
\begin{equation*}
p_1=p_{\mathrm{a}}\Big(1-\frac{11000}{44300}\Big)^{5.256}=0.223\,p_{\mathrm{a}}
\tag{3.66}
\end{equation*}

(ii) 同温层中压强分布

直接利用式(3.63)便可求出同温层中大气压强分布
\[
\frac{p}{p_1}=\exp\bigg[\frac{g(z_1-z)}{RT_1}\bigg]
\]
将$z_1=11000\,$m，$T_1=273-56.5=216.5\,$(K)，$p_1=0.223\,p_{\mathrm{a}}$代入后得
\begin{equation*}
p=0.223\,p_{\mathrm{a}}\exp[0.0001578(11000-z)]\,,\qquad z>11000\,\mathrm{m}
\tag{3.67}
\end{equation*}
式(3.65)和式(3.67)是国际标准大气的计算公式.

\subsection*{4. 重力场中静止流体作用在周围物体上的力}

水下结构物上受到的静水作用力，以及浮体上流体作用力和力矩等，都是工程设计所必需的数据.下面介绍静水作用力的计算方法及其主要性质.

(1)静水作用力公式

已知重力场中静止压强分布为(习惯上取垂直坐标与重力加速度方向一致，即$z$方向坐标向下为正)
\begin{equation*}
p=p_{\mathrm{a}}+\rho g z
\tag{3.68}
\end{equation*}
任意面上静水压强的合力$\boldsymbol{R}$为
\begin{equation*}
\boldsymbol{R}=\iint_A -p\boldsymbol{n}\,\mathrm{d}A
\tag{3.69}
\end{equation*}
合力矩为
\begin{equation*}
\boldsymbol{L}=\iint_A -p\boldsymbol{r}\times\boldsymbol{n}\,\mathrm{d}A
\tag{3.70}
\end{equation*}
式中$\boldsymbol{n}$为物体表面的法向量，它指向流场.我们需要强调指出：一般情况下合力$\boldsymbol{R}$与合力矩$\boldsymbol{L}$不垂直，或$\boldsymbol{R}\cdot\boldsymbol{L}\neq 0$，因此任意曲面上的静水作用力一般不能简化为等价力系.在某些特殊情况下，如平面上的静水作用力是平行力系，这时有$\boldsymbol{R}\perp\boldsymbol{L}$，可以把它简化为一等价力系，即可以简化为一个合力$\boldsymbol{R}$，它的作用线通过点$\boldsymbol{r}_0$，并由下式求出$\boldsymbol{r}_0$：
\begin{equation*}
\boldsymbol{r}_0\times\boldsymbol{R}=\boldsymbol{L}
\tag{3.71}
\end{equation*}
下面介绍几个计算静水作用力的实例.

(2)斜面上的静水作用力

设受力面与水平面($xy$平面)夹角等于$\alpha$，斜面法向量$\boldsymbol{n}$指向液体，如图3.15所示，如受力面另一侧受大气压强作用，这时受力面上静水压强的合力为
\[
\boldsymbol{R}=-\iint_A(p-p_{\mathrm{a}})\boldsymbol{n}\,\mathrm{d}A
\]
在重力场中静水压强$p=p_{\mathrm{a}}+\rho gz$，斜面的法向量$\boldsymbol{n}=\boldsymbol{j}\sin\alpha-\boldsymbol{k}\cos\alpha$为常向量，故合力向量
\[
\boldsymbol{R}=-\boldsymbol{n}\iint_A \rho gz\,\mathrm{d}A=-\rho g\boldsymbol{n}\iint_A z\,\mathrm{d}A
\]
根据面积中心(简称面心)的几何定义，它的坐标为$\boldsymbol{r}_c=\displaystyle\iint_A \boldsymbol{r}\,\mathrm{d}A/A$，故合力公式为
\begin{equation*}
\boldsymbol{R}=-\rho g z_c A\boldsymbol{n}
\tag{3.72}
\end{equation*}
$z_c$是受力面的面心在水面下的深度，于是，式(3.72)可表述为

\textbf{作用在水下受力斜面$A$上的静水作用力等于以该面积为底面，以该面积面心的深度为高的柱体中的水重，方向垂直于斜面.}

必须注意：式(3.72)并不表示静水作用力的作用点通过面心，合力作用点必须通过等价力系方法求出.下面先求静水中的合力矩.
\[
\boldsymbol{L}=-\iint_A(p-p_{\mathrm{a}})\boldsymbol{r}\times\boldsymbol{n}\,\mathrm{d}A=\boldsymbol{n}\times\iint_A \rho gz\boldsymbol{r}\,\mathrm{d}A
\]
为了便于计算上述积分，在受力平面上设立坐标系$x'y'$，其中$x'=x$，$y'=z/\sin\alpha$，见图3.15，这时被积函数中$\boldsymbol{r}=x'\boldsymbol{i}'+y'\boldsymbol{j}'$，力矩
\[
\boldsymbol{L}=\boldsymbol{n}\times\iint_A \rho g\,y'(\boldsymbol{x}'\boldsymbol{i}'+y'\boldsymbol{j}')\sin\alpha\,\mathrm{d}A=\rho g\boldsymbol{n}\times\Big(\boldsymbol{i}'\iint_A x'y'\,\mathrm{d}A+\boldsymbol{j}'\iint_A y'^{2}\,\mathrm{d}A\Big)\sin\alpha
\]
这里可以看到合力矩垂直于受力面的法向量，因而合力矩垂直于合力$\boldsymbol{R}$，所以作用在斜面上的静水作用力可以找到等价力系，设合力$\boldsymbol{R}$作用线通过点$\boldsymbol{r}_0$，它由以下公式确定：
\[
\boldsymbol{r}_0\times\boldsymbol{R}=\boldsymbol{L}
\]
将$\boldsymbol{R}$,$\boldsymbol{L}$的计算结果分别代入上式：
\[
-\rho g z_c A\,\boldsymbol{r}_0\times\boldsymbol{n}=\rho g\boldsymbol{n}\times\bigg(\iint_A x'y'\boldsymbol{i}'\,\mathrm{d}A+\iint_A y'^{2}\boldsymbol{j}'\,\mathrm{d}A\bigg)\sin\alpha
\]
利用向量代数运算(用$\boldsymbol{n}$叉乘上式)，可将上式简化为
\[
\boldsymbol{r}_0=\bigg(\frac{\boldsymbol{i}'\iint_A x'y'\,\mathrm{d}A}{z_c A}+\frac{\boldsymbol{j}'\iint_A y'^{2}\,\mathrm{d}A}{z_c A}\bigg)\sin\alpha
\]
由坐标变换$z_c=y_c'\sin\alpha$，得
\begin{equation*}
\boldsymbol{r}_0=\frac{\boldsymbol{i}'\iint_A x'y'\,\mathrm{d}A}{y_c'A}+\frac{\boldsymbol{j}'\iint_A y'^{2}\,\mathrm{d}A}{y_c'A}
\tag{3.73}
\end{equation*}
即合力作用点在受力面上的坐标为
\begin{equation*}
x_0'=\frac{\displaystyle\iint_A x'y'\,\mathrm{d}A}{y_c'A}\,,\qquad y_0'=\frac{\displaystyle\iint_A y'^{2}\,\mathrm{d}A}{y_c'A}
\tag{3.74}
\end{equation*}
式中$y_c'$是受力面面心的坐标.

\noindent\textbf{例3.12}\quad 圆心在水面下$y_c$处的垂直圆形闸门，半径为$a\,(a<y_c)$，计算圆形闸门上的静水作用力与合力作用点(见图3.16).

\noindent\textbf{解}\quad 圆形闸门上的作用力由式(3.72)得
\[
R=\pi\rho g\,y_c\,a^{2}
\]
作用点按式(3.74)计算，因圆形是对称图形，$\displaystyle\int x'y'\,\mathrm{d}A=0$，故$x_0'=0$；另一方面，由面积惯性矩公式$\displaystyle\iint_A y'^{2}\,\mathrm{d}A=\frac{\pi a^{4}}{4}+y_c^{2}A$，得
\[
y_0'=y_c+\frac{a^{2}}{4y_c}
\]
该例结果表明合力作用点在圆心以下$\dfrac{a^{2}}{4y_c}$处.

(3)封闭物体上静水作用力的合力：阿基米德(Archimedes)原理.

潜入水下封闭物体上静水作用的合力有简便计算公式
\begin{equation*}
\boldsymbol{R}=-\rho V\boldsymbol{g}
\tag{3.75}
\end{equation*}
并且合力通过潜水物体$V$的体心，$V$为物体排开水的体积.这一结果早在二千多年前由阿基米德发现，称为阿基米德原理.下面用静水作用力公式给予证明.计算合力及合力矩公式为
\[
\boldsymbol{R}=-\oint_\Sigma p\boldsymbol{n}\,\mathrm{d}A\,,\qquad \boldsymbol{L}=-\oint_\Sigma p\boldsymbol{r}\times\boldsymbol{n}\,\mathrm{d}A
\]
利用高斯公式，将两个面积分化为体积分如下：
\[
\boldsymbol{R}=\int_V -\nabla p\,\mathrm{d}V\,,\qquad \boldsymbol{L}=-\int_V \nabla\times(p\boldsymbol{r})\,\mathrm{d}V
\]
由静力学方程$\nabla p=\rho\boldsymbol{g}$，有
\[
\boldsymbol{R}=-\int_V \rho\boldsymbol{g}\,\mathrm{d}V=-\rho\boldsymbol{g}V
\]
力矩式中$\nabla\times(p\boldsymbol{r})=\nabla p\times\boldsymbol{r}+p\nabla\times\boldsymbol{r}=\rho\boldsymbol{g}\times\boldsymbol{r}$，故
\[
\boldsymbol{L}=-\int_V \rho\boldsymbol{g}\times\boldsymbol{r}\,\mathrm{d}V=-\rho\boldsymbol{g}\times\int_V \boldsymbol{r}\,\mathrm{d}V
\]
$\displaystyle\int_V \boldsymbol{r}\,\mathrm{d}V$可用体心公式表示，即$\displaystyle\int_V \boldsymbol{r}\,\mathrm{d}V=\boldsymbol{r}_c V$，最后得
\[
\boldsymbol{L}=-\rho\boldsymbol{g}V\times\boldsymbol{r}_c=\boldsymbol{R}\times\boldsymbol{r}_c
\]
该式说明$\boldsymbol{L}\perp\boldsymbol{R}$.用等价力系的概念，$(\boldsymbol{R},\boldsymbol{L})$力系等价于通过体心$\boldsymbol{r}=\boldsymbol{r}_c$的浮力$\boldsymbol{R}=-\rho V\boldsymbol{g}$.以上公式可表述为：\textbf{潜体上的静水作用力的大小等于潜体排开的水重，方向和重力加速度相反，并通过潜体的体心.}

(4)半潜体上浮力及浮力的平衡

物体的一部分浸入水中称为半潜体，半潜体上静水作用的合力仍应用式(3.75)计算，但式中$V$为排水体积，即水平面以下物体排开水的体积(图3.17).请读者自己证明这一结论.

\subsection*{5. 非惯性坐标系中静止流场的压强分布}

计算相对于运动坐标系静止的流场中压强分布时，应当在质量力场中包括惯性力，其他分析和计算方法与前面讨论的相同.在非惯性坐标系中，单位质量的惯性力等于$-\boldsymbol{a}$，因此非惯性坐标系中的静止平衡方程为
\begin{equation*}
\nabla p=\rho\boldsymbol{f}-\rho\boldsymbol{a}
\tag{3.76}
\end{equation*}
$\boldsymbol{a}$由式(3.17)给出.下面举两个例子来说明运动系统中静止流场压强分布的计算.

(1)匀加速度系统中静止液体在重力作用下的压强分布

体积力强度为$\rho\boldsymbol{g}-\rho\boldsymbol{a}$，因为匀加速度系统中$\boldsymbol{g}$,$\boldsymbol{a}$都是常数，因而质量力是有势的，它们的势很容易求出：
\[
\Pi=-\int\boldsymbol{f}\cdot\mathrm{d}\boldsymbol{x}=-\int(\boldsymbol{g}-\boldsymbol{a})\cdot\mathrm{d}\boldsymbol{x}=gx_3+a_1x_1+a_2x_2+a_3x_3
\]
令$(x_1,x_2,x_3)=(x,y,z)$($z$垂直向上)，并设系统在水平方向作匀加速度直线运动，即$\boldsymbol{a}=a\boldsymbol{i}$，这时相对于运动坐标系的外力场势为
\[
\Pi=gz+ax
\]
当$\rho=\mathrm{const.}$时压强场积分为$p/\rho+\Pi=C$，因此
\[
p=-\rho gz-\rho ax+C
\]
上式说明，在匀加速直线运动的坐标系中，等势面和等压面都不再是水平面，而是倾斜的平面.在运动系统中设立坐标轴，使它们的原点在自由面上，当已知自由面上压强$p_{\mathrm{a}}$时，液体中压强可由下式求出：
\begin{equation*}
p=-\rho gz-\rho ax+p_{\mathrm{a}}
\tag{3.77}
\end{equation*}

\noindent\textbf{例3.13}\quad 有一储液车(图3.18)，车身长$L$，当它作匀加速直线运动时，求储液罐前后液位的高度差.

\noindent\textbf{解}\quad 在储液车上设立坐标系$(x,z)$，加速度$\boldsymbol{a}$的正方向与$x$轴一致，同时坐标系原点在自由面上.根据本节所得等压面方程：
\[
p=-\rho gz-\rho ax+p_{\mathrm{a}}
\]
自由面方程为$p_{\mathrm{a}}=-\rho gz-\rho ax+p_{\mathrm{a}}$，即$gz+ax=0$，即储液车中自由液面的斜率为
\[
\frac{\mathrm{d}z}{\mathrm{d}x}=-\frac{a}{g}\qquad\text{或}\qquad \Delta z=-\frac{a}{g}\,\Delta x
\]
故前后端液位差：
\[
H=\frac{a}{g}\,L
\]
上式表明车厢前面的液位下降，车厢后面的液位上升，液面的高度差和储液车加速度成正比.

(2)在定轴等速旋转系统中的静止流场

如有一非惯性系统以重力加速度方向为定轴，作等角速度旋转运动.在该系统中静止的液体受到的体积力有重力$\rho\boldsymbol{g}$和惯性力$\rho\omega^{2}\boldsymbol{r}$，其中$\omega$为角速度，$\boldsymbol{r}$为柱坐标的向径(图3.19).

这时静平衡方程为
\begin{equation*}
\nabla p=\rho\boldsymbol{g}+\rho\omega^{2}\boldsymbol{r}
\tag{3.78}
\end{equation*}
体积力场$\rho\boldsymbol{g}+\rho\omega^{2}\boldsymbol{r}$是有势的，它的势为
\[
\Pi=-\int(\boldsymbol{g}+\omega^{2}\boldsymbol{r})\cdot\mathrm{d}\boldsymbol{x}=gz-\frac{1}{2}\omega^{2}r^{2}
\]
式中$\boldsymbol{r}$是柱坐标的向径，积分式(3.78)得压强分布：
\[
p=-\rho gz+\frac{1}{2}\omega^{2}r^{2}+C
\]
将柱坐标的原点$(z=0,r=0)$设立在自由面上，这时积分常数$C=p_{\mathrm{a}}$，故等角速度旋转系统的静止流场中压强分布为
\begin{equation*}
p=-\rho gz+\frac{1}{2}\omega^{2}r^{2}+p_{\mathrm{a}}
\tag{3.79}
\end{equation*}
自由液面$(p=p_{\mathrm{a}})$的方程为
\[
gz-\frac{1}{2}\omega^{2}r^{2}=0
\]
或
\begin{equation*}
z=\frac{\omega^{2}r^{2}}{2g}
\tag{3.80}
\end{equation*}
可见在等角速度定轴旋转系统中静止液体的自由表面(3.80)和等压面(3.79)均为旋转抛物面.
\section*{习题}\addcontentsline{toc}{section}{习题}

\noindent\textbf{3.1}\quad 图3.20所示一文丘利管，上游水箱的水位到喉部的高度为$h$，水箱横截面积为$S_0$.在文丘利管喉部有一小管与它下方的一水槽相连，水槽水位与喉部间高度为$L$，喉部截面积为$S_1$，文丘利管出口截面积为$S_2$.假定通过管道的流体是理想不可压缩的，$h$恒定不变，且$S_1/S_0\ll 1$，$S_2/S_0\ll 1$，重力加速度为$g$，水的密度为$\rho$，问$h$应至少为多高，方能将水槽中的水吸入文丘利管？

\medskip
\noindent\textbf{3.2}\quad 图3.21所示喷雾器，活塞以速度$U$作等速运动，喉部处空气造成低压，将液体吸入向大气喷雾.若空气密度为$\rho$，液体密度为$\rho'$，假定流动为不可压理想定常的，求能喷雾的最大吸入高度.

\medskip
\noindent\textbf{3.3}\quad 图3.22所示密度为$\rho$的不可压流体在水平变截面管内作一维定常流动，在管子位置$A,B$处连接一个U形管压力计.已知$A$处截面积为$F$，$B$处截面积为$f$，U形管内液体密度为$\rho'$.今测得U形管内两液面高度差为$h$，求管内通过的体积流量.

\medskip
\noindent\textbf{3.4}\quad 图3.23所示水平变截面管道，截面1-1处接引射管，截面2-2处通大气，截面积分别为$A_1$和$A_2$，当管内流过密度为$\rho$，流量为$Q$的不可压流体时，把密度为$\rho'$的流体吸入管道.求管道内能吸入$\rho'$流体的最大吸入管高度$h$，假定流动是一维定常的.

\medskip
\noindent\textbf{3.5}\quad 图3.24所示虹吸管吸水，假定流动是理想定常的，已知水的汽化压力为$2\,$kPa，大气压力为$101\,$kPa，试设计一个虹吸管，要求吸水量为$0.08\,\mathrm{m^3/s}$，管径为$10\,$cm，要求管内不出现气泡，求$h_1$和$h_2$.

\medskip
\noindent\textbf{3.6}\quad 如图3.25所示容器中盛水，高度为$h$，当打开阀门$B$时，求：

(1)出口水流速度随时间的变化；

(2)管道中压力分布；

(3)何时出口速度最大？最大值为多少？

(4)$L\to 0$时出口速度为多少？

假定容器的截面积远大于管道的截面积，流动过程中$h$保持不变，$p_{\mathrm{a}}$为大气压力.

\medskip
\noindent\textbf{3.7}\quad 图3.26所示等截面V形管内盛液体，液柱长度为$L$，当V形管在纸平面内作微小晃动后，求液面自由振荡的周期.

\medskip
\noindent\textbf{3.8}\quad 如图3.27所示，空气从$A,B$入口进入箱子，从$C$处流出.已知流动定常，$A,B$入口面积均为$5\,\mathrm{cm^2}$，出口面积为$10\,\mathrm{cm^2}$，$p_A=p_B=1.08\times 10^{5}\,\mathrm{N/m^2}$，$U_A=U_B=30\,$m/s，出口压力为大气压强$p_{\mathrm{a}}=1.03\times 10^{5}\,\mathrm{N/m^2}$，空气密度$\rho=1.23\,\mathrm{kg/m^3}$，求支撑力$F_1$和$F_2$.

\medskip
\noindent\textbf{3.9}\quad 水力采煤是用水枪在高压下喷出强力水柱冲击煤层(图3.28).设水枪出口直径为$30\,$mm，出口水流速度$U_1$为$50\,$m/s，求水柱对煤层的附加冲击力.假定流动是定常的，质量力和摩擦力可以忽略，水的密度为$1000\,\mathrm{kg/m^3}$.

\medskip
\noindent\textbf{3.10}\quad 喷气发动机作地面实验(图3.29).若质量流量为$\dot{m}_i$，喷气压力为$p_i$，喷气速度为$U_i$，喷管出口面积为$A$，大气压强为$p_{\mathrm{a}}$，流动是定常的，求试验台所受的推力$F$.

\medskip
\noindent\textbf{3.11}\quad 为了测定圆柱体的阻力系数$C_D$，将一个直径为$d$，长度为$l$的圆柱放在二维定常不可压缩流中，实验在风洞中进行，在图3.30中1-1、2-2截面上测得近似的速度分布如图.这两个截面上的压力都是均匀的，等于$p_\infty$.试求圆柱体的阻力系数$C_D$.$C_D$的定义为
\[
C_D=\frac{2D}{\rho U_\infty^{2}\,ld}
\]
其中$D$为圆柱绕流时的阻力(流体对圆柱的作用力的水平方向分量)，$\rho$为流体密度，$U_\infty$为来流速度.

\medskip
\noindent\textbf{3.12}\quad 证明不可压缩流体通过突然扩大的管道时压力损失$\Delta h=U_1^{2}\,(1-A_1/A_2)^{2}/2g$(图3.31)，这里压头损失的定义是$\Delta h=(p_1/\gamma+U_1^{2}/2g)-(p_2/\gamma+U_2^{2}/2g)$.假定壁面上没有摩擦力，1-1截面上压力均匀并等于$p_1$，2-2截面上压力与速度也是均匀的，流动是定常的.

\medskip
\noindent\textbf{3.13}\quad 如图3.32所示洒水器，两边喷水，同时绕铅垂轴作等速转动.若喷嘴直径为$12\,$mm，每个喷嘴的喷水量$Q$为$80\,\mathrm{cm^3/s}$，忽略机械摩擦，旋转轴上没有任何扭矩，求洒水器旋转角速度$\omega$.

\medskip
\noindent\textbf{3.14}\quad 图3.33所示风机叶轮的内径$d_1=12.5\,$cm，外径$d_2=30\,$cm，叶片宽度$b=2.5\,$cm，转速$n=1725\,$r/min，体积流量$Q=372\,\mathrm{m^3/h}$，空气沿径向流入叶轮$(\alpha_1=90^{\circ})$.进口压强$p_1=9.7\times 10^{4}\,\mathrm{N/m^2}$，气温$t_1=20\,^{\circ}$C，气体流出气角$\beta_2=30^{\circ}$(假定流体是理想不可压缩的).

(1)计算$c_1,u_1,w_1,\beta_1$及$c_2,u_2,w_2,\alpha_2$；

(2)求所需的扭矩.

\medskip
\noindent\textbf{3.15}\quad 空气以$U_1=100\,$m/s进入气罐，又以$U_2=200\,$m/s离开气罐，如果流动是绝热的，空气也不做功，证明出口处空气温度比进口处低$14.9\,$K.

\medskip
\noindent\textbf{3.16}\quad 有一射流泵(引射器)如图3.34所示.高速流体从主流道中在1-1截面上以速度$U_1$喷出，带动流道中的流体在1-1截面上以速度$U_2$流出.假定二者为同一不可压流体，这两股流体在等直径的混合室中由于流体间的摩擦和相互掺混，在出口截面3-3处速度变为均匀.假定1-1截面上两股流体压力相同，混合室表面的摩擦力可以忽略，流道是定常的.现测得$A_1=0.0093\,\mathrm{m^2}$，$A_3=0.093\,\mathrm{m^2}$，$U_1=30.48\,$m/s，$U_2=3.04\,$m/s，$\rho=1000\,\mathrm{kg/m^3}$，试求$U_3$和$p_3-p_1$值.这里$A_1$是主流出口面积.

\medskip
\noindent\textbf{3.17}\quad 在管系中有一$90^{\circ}$弯管放在水平面上(图3.35)，管径$d_1=15\,$cm，$d_2=7.5\,$cm，入口处水的平均流速$U_1=25\,$m/s，压强$p_1=6.86\,\mathrm{N/cm^2}$(表压)，不计阻力损失，流动为定常，试求支撑弯管所需的力.

\medskip
\noindent\textbf{3.18}\quad 图3.36所示斜平板迎着射流以等速$U_0$运动，$A_0$和$V_0$分别为固定喷口的出口面积和水流的绝对速度.不计重力和摩擦力，并假定相对于斜平板的流动是定常的，水的密度为$\rho$.试导出维持平板运动所需功率.

\medskip
\noindent\textbf{3.19}\quad 流量为$Q_0$的自由射流与平板相遇(图3.37)，其中一股流量为$Q_1$的流动偏转角为$90^{\circ}$，其余流动偏转角为$\alpha$，如图所示.已知射流速度$U_0=20\,$m/s，总流量$Q_0=24\times 10^{-3}\,\mathrm{m^3/s}$及$Q_1=8\times 10^{-3}\,\mathrm{m^3/s}$，$\rho=1000\,\mathrm{kg/m^3}$，不计水的重量和粘性，并假定流动定常，在足够远处$U_1,U_2$均匀.求：(1)$Q_2,U_1,U_2,\alpha$；(2)平板所受的力.

\medskip
\noindent\textbf{3.20}\quad 如图3.38所示，密度为$\rho$的不可压流体在$A$-$A$截面处均匀进入$AB$段圆管，由于粘性作用，在$B$-$B$截面处以抛物线速度分布流出.测得进口$A$-$A$截面上压强为$p$，速度为$U$，大气压强为$p_{\mathrm{a}}$.求管内流动作用在法兰上的力$F$.

\medskip
\noindent\textbf{3.21}\quad 如图3.39所示，水箱1中的水经光滑无阻力的圆孔口水平射出，冲到一平板上.平板封盖着另一水箱2的孔口，水箱1中水位高度为$h_1$，水箱2中水位高度为$h_2$，两孔口中心重合，而且$d_1=d_2/2$.若射流的形状是对称的，冲击到平板后，转向平行于平板的方向，并向四周均匀流出.假定流动为理想不可压定常流动，平板摩擦阻力和质量力不计.当$h_1$为已知时，问$h_2$最大不能超过多少？
\medskip
\noindent\textbf{3.22}\quad 已知
\[
\boldsymbol{P}=\begin{bmatrix}
p_{xx} & p_{xy} & p_{xz}\\
p_{yx} & p_{yy} & p_{yz}\\
p_{zx} & p_{zy} & p_{zz}
\end{bmatrix}
\]
及坐标变换的方向余弦
\[
\begin{array}{c|ccc}
 & i & j & k\\ \hline
i' & l_1 & m_1 & n_1\\
j' & l_2 & m_2 & n_2\\
k' & l_3 & m_3 & n_3
\end{array}
\]
$l_i,m_i,n_i$为坐标轴间的方向余弦，求$p_{x'x'}$和$p_{y'z'}$的表达式.

\medskip
\noindent\textbf{3.23}\quad 证明正交曲线坐标系中的$\nabla\cdot(\boldsymbol{P}\cdot\boldsymbol{U})$的表达式为
\[
\nabla\cdot(\boldsymbol{P}\cdot\boldsymbol{U})=\delta_{ik}\,\frac{1}{h_k}\,\frac{\partial}{\partial x_k}(p_{ij}U_j)+\frac{p_{ij}U_j}{h_ih_k}\,\frac{\partial h_k}{\partial q_i}\bigg|_{k\neq i}
\]
$\boldsymbol{P}$是应力张量，$p_{ij}$是应力分量.

\medskip
\noindent\textbf{3.24}\quad 写出下列流体运动的连续性方程：

(1)流体质点以角速度$\omega$绕$z$轴作圆周运动；

(2)流体轨迹位于绕$z$轴的圆柱面上；

(3)流体质点在包含$z$轴的平面上运动；

(4)流体质点位于与$z$轴同轴并有共同顶点(原点)的圆锥上.

\medskip
\noindent\textbf{3.25}\quad 写出下列流体运动的连续性方程：

(1)流体质点作径向运动，且$\boldsymbol{U}=U(R,t)\boldsymbol{e}_R$；

(2)流体质点在同心球面上运动.

\medskip
\noindent\textbf{3.26}\quad 证明：在$\omega=\mathrm{const.}$的坐标系中，当理想气体均熵流动，并忽略质量力时，有
\[
\frac{D'}{Dt}\bigg(\frac{U'^{2}}{2}+h-\frac{U_{\mathrm{c}}^{2}}{2}\bigg)=\frac{1}{\rho}\,\frac{\partial' p}{\partial t}
\]
式中$h=e+\dfrac{p}{\rho}$为单位质量的焓，$U_{\mathrm{c}}$为牵连速度的模，上角标“$'$”为相对于动坐标系中的导数或物理量.

\medskip
\noindent\textbf{3.27}\quad 水流通过截面积为$s\times s$的正方形圆环形通道(图3.40)，其中心线的曲率半径为$R$，测得通道外侧与内侧压差为$\Delta p$，假定流动为理想定常的，水流密度为$\rho$，且已知通道内速度分布$U_z=U_r=0$，$U_\theta=A/r$，这里$A$为待定常数.试证明通道内流量为
\[
Q=\Big(R^{2}-\frac{s^{2}}{4}\Big)\sqrt{\frac{s\Delta p}{R\rho}}\,\ln\frac{2R+s}{2R-s}
\]

\medskip
\noindent\textbf{3.28}\quad 证明静止的非正压流场中质量力必满足$\boldsymbol{f}\cdot\nabla\times\boldsymbol{f}=0$，又问若$\boldsymbol{f}\cdot\nabla\times\boldsymbol{f}=0$，流场是否一定静止.

\medskip
\noindent\textbf{3.29}\quad 给出如下质量力场，分别在(a)正压或不可压流场；(b)斜压流场$\rho\neq\rho(p)$中说明流场是否可能静止.

(1)$\boldsymbol{f}=(y^{2}+yz+z^{2})\boldsymbol{i}+(z^{2}+zx+x^{2})\boldsymbol{j}+(x^{2}+xy+y^{2})\boldsymbol{k}$；

(2)$\boldsymbol{f}=-\dfrac{k}{R^{2}}\boldsymbol{e}_R$，$k$为常数，$R>0$(球坐标).

\medskip
\noindent\textbf{3.30}\quad 一倾斜式微压计，由一储液杯和一斜管组成，杯内自由液面上的压力为$p_{\mathrm{a}}$(图3.41).斜管一端接待测压容器，设待测压力为$p_1$，液体重度$\gamma=\rho g$，若$p_1=p_{\mathrm{a}}$时，斜管中液柱读数为$a_0$，试证明：
\[
p_{\mathrm{a}}-p_1=\gamma(a-a_0)\bigg(1+\frac{A_1}{A_2}\,\frac{1}{\sin\alpha}\bigg)\sin\alpha=K(a-a_0)
\]
式中$A_1$为斜管截面积，$A_2$为储液杯截面积，$\alpha$为斜管倾斜角，$a$为测压时斜管液柱的读数，$K=\gamma\Big(1+\dfrac{A_1}{A_2}\dfrac{1}{\sin\alpha}\Big)\sin\alpha$.又若已知$K=4\,\mathrm{kN/m^3}$，$a-a_0=65\,$mm，求$p_{\mathrm{a}}-p_1$.

\medskip
\noindent\textbf{3.31}\quad 有一差式测压管，连通方式如图3.42所示，若测得$a,b,c$值，求1-1截面和2-2截面上的压力差$p_1-p_2$.设连通管中两种液体的密度$\rho$和$\rho'$为已知.

\medskip
\noindent\textbf{3.32}\quad 直径为$1.5\,$m气象用的探空气球，携带一仪器盒，质量为$0.4\,$kg.仪器盒体积可以略去不计，求此气球可能上升的高度.

\medskip
\noindent\textbf{3.33}\quad 如图3.43所示水坝的平板闸门，求开启时所需的提升力$F$.已知闸门与闸墩之间的摩擦系数为$0.2$，闸门质量为$10000\,$kg，闸门高$3\,$m，宽$10\,$m.

\medskip
\noindent\textbf{3.34}\quad 如图3.44所示水坝，求开启弧形闸门时所需的提升力$F$.已知弧形闸门的圆弧半径为$9\,$m，闸门宽度为$10\,$m，闸门和桁架的重量为$1176\times 10^{3}\,$N，重心与$O$点的水平距离为$7\,$m，闸门的作用力距$O$点为$10\,$m，不计闸门和闸墩之间的摩擦力.图中$A$点为圆弧几何中心，$O$点为闸门铰链中心.

\medskip
\noindent\textbf{3.35}\quad 如图3.45所示，求流体垂直作用于半径为$a$的四分之一圆板上的总作用力$F$和压力中心$C$的位置.已知$Ox$与流体自由水平面重合，自由面上压力为零.

\medskip
\noindent\textbf{3.36}\quad 半径为$a$具有垂直轴$Oz$的半球内盛满液体(图3.46).求被$Oxz$和$Oyz$两正交平面所切开的四分之一球面上压力的合力$F$及压力中心$c$的位置(合力作用于球面的位置).水面上压力为零.

\medskip
\noindent\textbf{3.37}\quad 一个充满水的密闭容器，以等角速度$\omega$绕一水平轴旋转.证明它的等压面为圆柱面，且该圆柱面的轴线比转动轴高$g/\omega^{2}$.

\medskip
\noindent\textbf{3.38}\quad (儒可夫斯基疑题)在一盛满液体的容器的垂直壁面上装一均匀圆柱(图3.47)，可以无摩擦地绕水平轴$O$旋转，圆柱的一半在所有时间均保持沉没于液体之内.根据阿基米德原理，似乎圆柱受到一个向上的力可以迫使圆柱旋转，这样似乎不要消耗能量便可以得到功，即永动机可以实现了.试说明为何圆柱不会旋转？求液体作用于圆柱的合力$F$及其压力中心的位置(合力作用于圆柱表面的位置).

\medskip
\noindent\textbf{3.39}\quad 利用装有液体并与物体一起运动的U形管测量物体的加速度，如图3.48所示.U形管直径很小，$l=30\,$cm，$h=5\,$cm.求物体加速度$a$.

\medskip
\noindent\textbf{3.40}\quad 如图3.49所示，密闭容器内的气体压强$p_b$(小于大气压强)始终保持不变，U形管随容器一起运动.试建立U形管读数$h$和$a,\omega,r_1,r_2,p_b,\rho$之间的关系.设(1)容器垂直向上运动，加速度为$a$；(2)容器以角速度$\omega$绕中心轴旋转.

\medskip
\noindent\textbf{3.41}\quad 盛满水的容器长为$L$，高为$H$，宽为1，在右上角$C$处与大气相通，它以等加速度$a_0$沿倾斜角为$\alpha$的坡面向上行驶，求容器内$A$角及$B$角处的压力(图3.50).

\medskip
\noindent\textbf{3.42}\quad 盛水圆筒绕一铅垂轴以$\omega$等速转动(图3.51).圆筒半径为$a$，圆筒中心$O'$到转动轴的距离为$R$，静止水面高度为$h$，求自由面形状.

\section*{思考题}\addcontentsline{toc}{section}{思考题}

\noindent\textbf{S3.1}\quad 船上安装水泵从船外吸入河水，经水泵再喷出船外，这种装置是否能推动该船前进？为什么？

\medskip
\noindent\textbf{S3.2}\quad 由螺旋桨推进的飞机，桨前方的空气速度(相对于飞机)和桨后的速度(相对于飞机)应有什么关系？

\medskip
\noindent\textbf{S3.3}\quad 应用理想不可压缩流体模型求解两股平行但流速不同的汇合流动，或从直管流入突然扩张/突然收缩管道时，流体的动量和动能是否有损失？

\medskip
\noindent\textbf{S3.4}\quad 在旋转系统中研究流体运动，除了重力以外，还有什么体积力？这些体积力是否对流体做功？在旋转系统中用文丘里管测量流速，其计算公式和静止系统是否相同？